% ============ Chapter 6. Regular sequences (原书 pp.123-153) ============
\chapter{Regular sequences}\label{ch:6}

In the 1950s homological algebra was introduced into commutative ring
theory, opening up new avenues of study. In this chapter we run through
some fundamental topics in this direction.

In \S\ref{sec:16} we define regular sequences, depth and the Koszul complex. The
notion of depth is not very geometric, and rather hard to grasp, but is
an extremely important invariant. It can be treated either in terms of
$\Ext$'s, or by means of the Koszul complex, and we give both versions. We
discuss the relation between regular and quasi-regular sequences in a
transparent treatment due to Rees. \S\ref{sec:17} contains the definition and principal
properties of Cohen--Macaulay (CM) rings. The theorem that quotients of
CM rings are always catenary is of great significance in dimension
theory. In \S\ref{sec:18} we treat a distinguished subclass of CM rings having even
nicer properties, the Gorenstein rings. In the famous paper of H.\ Bass~\cite{Ba1},
Gorenstein rings are discussed using Matlis' theory of injective modules.
Here we give an elementary treatment of Gorenstein rings following~\cite{G6}
before going through Matlis theory.

% ============ §16 Regular sequences and the Koszul complex (原书 pp.123-132) ============
\section{Regular sequences and the Koszul complex}\label{sec:16}\index{regular sequences, section 16}\index{Koszul complex, section 16}

Let $A$ be a ring and $M$ an $A$-module. An element $a \in A$ is said
to be $M$-regular\index{regular element} if $ax \neq 0$ for all $0 \neq x \in M$. A sequence $a_1, \ldots, a_n$ of elements
of $A$ is an \emph{$M$-sequence}\index{M-sequence} (or an \emph{$M$-regular sequence}) if the following two
conditions hold:
\begin{enumerate}[(1)]
  \item $a_1$ is $M$-regular, $a_2$ is $(M/a_1M)$-regular, $\ldots$, $a_n$ is $(M/\sum_{i=1}^{n-1} a_iM)$-
        regular;
  \item $M \neq \sum_{i=1}^{n} a_iM$.
\end{enumerate}
Note that, after permutation, the elements of an $M$-sequence may no longer
form an $M$-sequence.

\begin{theorem}\label{thm:16.1}
If $a_1, \ldots, a_n$ is an $M$-sequence then so is $a_1^{\nu_1}, \ldots, a_n^{\nu_n}$ for
positive integers $\nu_1, \ldots, \nu_n$.
\end{theorem}

\begin{proof}
It is sufficient to prove that if $a_1, \ldots, a_n$ is an $M$-sequence then so is
$a_1^{\nu_1}, a_2, \ldots, a_n$. Indeed, assuming this, we have in turn that $a_1^{\nu_1}, a_2, \ldots, a_n$
is an $M$-sequence, then setting $M_1 = M/a_1^{\nu_1}M$ we have that $a_2, a_3, \ldots, a_n$,
and hence also $a_2^{\nu_2}, a_3, \ldots, a_n$, is an $M_1$-sequence, and so on. Also, the
second condition $M \neq \sum a_i^{\nu_i}M$ is obvious.

Let us now prove by induction on $n$ that if $b_1, \ldots, b_n$ is an $M$-sequence,
and if $\sum_{i=1}^{n} b_it_i = 0$ with $t_i \in M$ then $t_i \in b_1M + \cdots + b_nM$ for all
$i$. First of all from the condition that $b_n$ is not a zero-divisor modulo
$b_1, \ldots, b_{n-1}$ we can write
\[
  t_n = \sum_{i=1}^{n-1} b_i\eta_i, \quad \text{with } \eta_i \in M.
\]
Therefore $\sum_{i=1}^{n-1} b_i(t_i + b_n\eta_i) = 0$, so that by induction we have
\[
  t_i + b_n\eta_i \in b_1M + \cdots + b_{n-1}M \quad \text{for } 1 \le i \le n - 1,
\]
giving $t_i \in b_1M + \cdots + b_nM$ for $1 \le i \le n - 1$. The condition for $t_n$ is already known.

Now assuming $\nu > 1$ we prove by induction on $\nu$ that $a_1^{\nu}, a_2, \ldots, a_n$ is
an $M$-sequence. Since $a_1$ is $M$-regular, so is $a_1^{\nu}$. For $i > 1$, suppose that
for some $\sigma \in M$ we have
\[
  a_i\sigma = a_1^{\nu}\xi_1 + a_2\xi_2 + \cdots + a_{i-1}\xi_{i-1} \quad \text{with } \xi_j \in M.
\]
Then since $a_1^{\nu-1}, a_2, \ldots, a_i$ is an $M$-sequence, we can write
\[
  \sigma = a_1^{\nu-1}\eta_1 + \cdots + a_{i-1}\eta_{i-1} \quad \text{with } \eta_j \in M.
\]
Hence we get
\[
  0 = a_1^{\nu-1}(a_i\eta_1 - a_i\xi_1) + a_2(\eta_2 - a_i\xi_2) + \cdots + a_{i-1}(\eta_{i-1} - a_i\xi_{i-1}).
\]
The above assertion gives $a_i\eta_1 - a_i\xi_1 \in a_1^{\nu-1}M + a_2M + \cdots + a_{i-1}M$,
and hence $a_i\eta_1 \in a_1M + a_2M + \cdots + a_{i-1}M$. Therefore $\eta_1 \in a_1M + \cdots +
a_{i-1}M$, and so as required we have $\sigma \in a_1^{\nu}M + a_2M + \cdots + a_{i-1}M$.
\end{proof}

Let $A$ be a ring, $X_1, \ldots, X_n$ indeterminates over $A$, and $M$ an $A$-module.
We can view elements of $M \otimes_A A[X_1, \ldots, X_n]$ as polynomials in the $X_i$
with coefficients in $M$,
\[
  F(X) = F(X_1, \ldots, X_n) = \sum t_{\alpha_1 \cdots \alpha_n}X_1^{\alpha_1} \cdots X_n^{\alpha_n}, \quad \text{with } t_{\alpha} \in M.
\]
For this reason we write $M[X_1, \ldots, X_n]$ for $M \otimes_A A[X_1, \ldots, X_n]$; we can
consider this either as an $A$-module or as an $A[X_1, \ldots, X_n]$-module. For
$a_1, \ldots, a_n \in A$ and $F \in M[X_1, \ldots, X_n]$, we can substitute the $a_i$ for $X_i$ to
get $F(a_1, \ldots, a_n) \in M$.

\begin{definition}
Let $a_1, \ldots, a_n \in A$, set $I = \sum a_iA$ and let $M$ be an $A$-module
with $IM \neq M$. We say that $a_1, \ldots, a_n$ is an $M$-quasi-regular sequence\index{quasi-regular sequence} if the
following condition holds for each $\nu$:
\[
  (*) \quad
  \begin{aligned}
    &F(X_1, \ldots, X_n) \in M[X_1, \ldots, X_n] \text{ is homogeneous of degree } \nu \text{ and }\\
    &F(a) \in I^{\nu+1}M \text{ implies that all the coefficients of } F \text{ are in } IM.
  \end{aligned}
\]
This notion is obviously independent of the order of $a_1, \ldots, a_n$.
\end{definition}

In the above definition it would not make any difference if we replaced
the condition that $F(a) \in I^{\nu+1}M$ by the condition $F(a) = 0$. Indeed, if $F$ is
homogeneous of degree $\nu$ and $F(a) \in I^{\nu+1}M$ then there exist a homogeneous
element $G(X) \in M[X_1, \ldots, X_n]$ of degree $\nu + 1$ such that $F(a) = G(a)$. Then
write $G(X) = \sum X_iG_i(X)$ with each $G_i$ homogeneous of degree $\nu$, and set
$F^*(X) = F(X) - \sum a_iG_i(X)$, so that $F^*$ is homogeneous of degree $\nu$ and
$F^*(a) = 0$. Moreover, if $F^*$ has coefficients in $IM$ then so does $F$.

We can define a map $\varphi\colon (M/IM)[X_1, \ldots, X_n] \to \gr_I M = \bigoplus_{\nu \ge 0}
I^\nu M/I^{\nu+1}M$ as follows: taking a homogeneous element $F(X) \in M[X]$ of
degree $\nu$ into the class of $F(a)$ in $I^\nu M/I^{\nu+1}M$ provides a homomorphism (of
additive groups) from $M[X]$ into $\gr_I M$ which preserves degrees. Since
$IM[X]$ is in the kernel, this induces a homomorphism
\[
  \varphi\colon M[X]/IM[X] = (M/IM)[X] \to \gr_I M,
\]
which is obviously surjective. Then $a_1, \ldots, a_n$ is a quasi-regular sequence
precisely when $\varphi$ is injective, and hence an isomorphism.

\begin{theorem}\label{thm:16.2}
Let $A$ be a ring, $M$ an $A$-module, and $a_1, \ldots, a_n \in A$; set
$I = (a_1, \ldots, a_n)A$. Then we have the following:
\begin{enumerate}[(i)]
  \item if $a_1, \ldots, a_n$ is an $M$-sequence then it is $M$-quasi-regular;
  \item if $a_1, \ldots, a_n$ is an $M$-quasi-regular sequence, and if $x \in A$ satisfies
        $IM:x = IM$ then $I^\nu M:x = I^\nu M$ for any $\nu > 0$.
\end{enumerate}
\end{theorem}

\begin{proof}[Proof (taken from Rees [5])]
First of all we prove (ii) by induction on $\nu$.
The case $\nu = 1$ is just the assumption; suppose that $\nu > 1$. For $\xi \in M$, if
$x\xi \in I^\nu M$ then also $x\xi \in I^{\nu-1}M$, so that by the inductive hypothesis
$\xi \in I^{\nu-1}M$, and hence we can write $\xi = F(a)$ with $F = F(X) \in M[X_1, \ldots,
X_n]$ homogeneous of degree $\nu - 1$. Now $x\xi = xF(a) \in I^\nu M$, so that by
definition of quasi-regular sequence each coefficient of $xF(X)$ belongs to
$IM$. Using $IM:x = IM$ once more we find that the coefficients of $F(X)$
also belong to $IM$, and therefore $\xi = F(a) \in I^\nu M$.

Now we prove (i) by induction on $n$. The case $n = 1$ can easily be
checked. Suppose that $n > 1$, and that the statement holds up to $n - 1$, so
that in particular $a_1, \ldots, a_{n-1}$ is $M$-quasi-regular. Now let $F(X) \in
M[X_1, \ldots, X_n]$ be homogeneous of degree $\nu$, such that $F(a) = 0$. We prove
by induction on $\nu$ that the coefficients of $F$ belong to $IM$. We separate
out $F(X)$ into terms containing $X_n$ and not containing $X_n$, writing
\[
  F(X) = G(X_1, \ldots, X_{n-1}) + X_nH(X_1, \ldots, X_n).
\]
Here $G$ is homogeneous of degree $\nu$ and $H$ of degree $\nu - 1$. Then, as we
proved in (ii),
\[
  H(a) \in (a_1, \ldots, a_{n-1})^\nu M:a_n = (a_1, \ldots, a_{n-1})^\nu M \subset I^\nu M,
\]
and hence, by induction on $\nu$, the coefficients of $H(X)$ belong to $IM$.
Moreover, by the above formula there is a homogeneous polynomial
$h(X_1, \ldots, X_{n-1})$ of degree $\nu$ with coefficients in $M$ such that $H(a) =
h(a_1, \ldots, a_{n-1})$, and so setting
\[
  g(X) = G(X_1, \ldots, X_{n-1}) + a_nh(X_1, \ldots, X_{n-1}),
\]
since $a_1, \ldots, a_{n-1}$ is $M$-quasi-regular, the coefficients of $g$ lie in
$(a_1, \ldots, a_{n-1})M$; therefore the coefficients of $G$ lie in
\[
  (a_1, \ldots, a_n)M.
\]
\end{proof}

This theorem holds for any $A$ and $M$, but as we will see in the next
theorem, under some conditions we can say that conversely, quasi-regular
implies regular. Then the notions of regular and quasi-regular sequences
for $M$ coincide, and so reordering an $M$-sequence gives again an $M$-
sequence.

\begin{theorem}\label{thm:16.3}
Let $A$ be a Noetherian ring, $M \neq 0$ an $A$-module, and
$a_1, \ldots, a_n \in A$; set $I = (a_1, \ldots, a_n)A$. Under the condition
\[
  (*) \quad \text{each of } M, M/a_1M, \ldots, M/(a_1, \ldots, a_{n-1})M \text{ is } I\text{-adically separated},
\]
if $a_1, \ldots, a_n$ is $M$-quasi-regular it is an $M$-sequence.
\end{theorem}

\begin{remark}
The hypothesis $(*)$ holds in either of the following cases:
\begin{enumerate}[($\beta$)]
  \item $M$ is finite and $I \subset \rad(A)$;
  \item $A$ is an $\N$-graded ring, $M$ an $\N$-graded module, and each $a_i$ is
        homogeneous of positive degree.
\end{enumerate}
However, for a non-Noetherian ring $A$ there are examples where the
theorem fails (Dieudonn\'e [1]) even if $A$ is local, $M = A$ and $I \subset \rad(A)$.
\end{remark}

\begin{proof}
We prove first that $a_1$ is $M$-regular. If $\xi \in M$ with $a_1\xi = 0$ then
by hypothesis $\xi \in IM$. Then setting $\xi = \sum a_i\eta_i$ we get $0 = \sum a_1a_i\eta_i$, so that
$\eta_i \in IM$. Proceeding in the same way we get $\xi \in \bigcap_{\nu \ge 1} I^\nu M = (0)$.

Now set $M_1 = M/a_1M$; if we prove that $a_2, \ldots, a_n$ is an $M_1$-quasi-regular
sequence then the theorem follows by induction on $n$. (If $M$ is $I$-adically
separated and $M \neq 0$ then $M \neq IM$.) So let $\bar{f}(X_2, \ldots, X_n)$ be a homogeneous
polynomial of degree $\nu$ with coefficients in $M_1$, such that $\bar{f}(a_2, \ldots, a_n) = 0$.
If $F(X_1, \ldots, X_n)$ is a homogeneous polynomial of degree $\nu$ with coefficients
in $M$ which reduces to $\bar{f}$ modulo $a_1M$, then $F(a_1, \ldots, a_n) \in a_1M$. Set
$F(a_1, \ldots, a_n) = a_1\sigma$; suppose that $\sigma \in I^iM$, so that we can write $\sigma = G_i(a)$
with $G_i(X) \in M[X_1, \ldots, X_n]$ homogeneous of degree $i$. Then
\[
  a_1G_i(a_2, \ldots, a_n) = a_1G_i(a_1, \ldots, a_n),
\]
and if $i < \nu - 1$ it follows that the coefficients of $G_i$ belong to $IM$, so that
$\sigma \in I^{i+1}M$; repeating this argument we see that $\sigma \in I^{\nu-1}M$. Setting $i = \nu - 1$
in the above formula, then since $X_1$ does not appear in $F$, we can
apply the definition of quasi-regular sequence to $F(X_1, \ldots, X_n) -
X_1G_{\nu-1}(X_1, \ldots, X_n)$ to deduce that the coefficients of $F$ belong to $IM$.
Hence, the coefficients of $\bar{f}$ belong to $IM_1$.
\end{proof}

\begin{corollary*}
Let $A$ be a Noetherian ring, $M$ an $A$-module and $a_1, \ldots, a_n$
an $M$-sequence. If conditions ($\alpha$) or ($\beta$) of the above remark hold then any
permutation of $a_1, \ldots, a_n$ is again an $M$-sequence.
\end{corollary*}

Here is an example where a permutation of an $M$-sequence fails to be
an $M$-sequence: let $k$ be a field, $A = k[X, Y, Z]$ and set $a_1 = X(Y - 1)$,
$a_2 = Y$, $a_3 = Z(Y - 1)$. Then $(a_1, a_2, a_3)A = (X, Y, Z)A \neq A$, and $a_1, a_2, a_3$
is an $A$-sequence, whereas $a_1, a_3, a_2$ is not.

\subsection*{The Koszul complex}

Given a ring $A$ and $x_1, \ldots, x_n \in A$, we define a complex $K_\bullet$ as follows:
set $K_0 = A$, and $K_p = 0$ if $p$ is not in the range $0 \le p \le n$. For $1 \le p \le n$,
let $K_p = \bigoplus A e_{i_1 \cdots i_p}$ be the free $A$-module of rank $\binom{n}{p}$ with basis
$\{e_{i_1, \ldots, i_p} \mid 1 \le i_1 < \cdots < i_p \le n\}$. The differential $d\colon K_p \to K_{p-1}$ is defined
by setting
\[
  d(e_{i_1 \cdots i_p}) = \sum_{r=1}^{p} (-1)^{r-1}x_{i_r}e_{i_1 \cdots \hat{i}_r \cdots i_p}
\]
(for $p = 1$, set $d(e_i) = x_i$). One checks easily that $dd = 0$. This complex is
called the \emph{Koszul complex}\index{Koszul complex}, and written $K_\bullet(x_1, \ldots, x_n)$ (alternatively,
$K_\bullet(x)$ or $K_{\bullet, x}$). For an $A$-module $M$ we set $K_\bullet(x, M) = K_\bullet(x) \otimes_A M$.
Moreover, for a complex $C_\bullet$ of $A$-modules we set $C_\bullet(x) = C_\bullet \otimes K_\bullet(x)$. In
particular, for $n = 1$ the complex $K_\bullet(x)$ is just
\[
  0 \to 0 \to A \xrightarrow{x} A \to 0,
\]
and it is easy to check that $K_\bullet(x_1, \ldots, x_n) = K_\bullet(x_1) \otimes \cdots \otimes K_\bullet(x_n)$. Since the
tensor product of complexes satisfies $L_\bullet \otimes M_\bullet \cong M_\bullet \otimes L_\bullet$, the Koszul com-
plex is invariant (up to isomorphism) under permutation of $x_1, \ldots, x_n$. The
Koszul complex $K_\bullet(x, M)$ has homology groups $H_p(K_\bullet(x, M))$, which we
abbreviate to $H_p(x, M)$. Quite generally we have
\[
  H_0(x, M) = M/\mathfrak{a}M,
\]
where $\mathfrak{a} = \sum x_iA$, and
\[
  H_n(x, M) \cong \{m \in M \mid x_1m = \cdots = x_nm = 0\}.
\]

\begin{theorem}\label{thm:16.4}
Let $C_\bullet$ be a complex of $A$-modules and $x \in A$. Then we obtain
an exact sequence of complexes
\[
  0 \to C_\bullet \to C_\bullet(x) \to C'_\bullet \to 0,
\]
where $C'_\bullet$ is the complex obtained by shifting the degrees in $C_\bullet$ up by $1$ (that is
$C'_p = C_{p-1}$, and the differential of $C'_\bullet$ is that of $C_\bullet$). The homology long
exact sequence obtained from this is
\[
  \cdots \to H_p(C_\bullet) \to H_p(C_\bullet(x)) \to H_{p-1}(C_\bullet) \xrightarrow{(-1)^{p-1}x}
  H_{p-1}(C_\bullet) \to \cdots;
\]
we have $x \cdot H_p(C_\bullet(x)) = 0$ for all $p$.
\end{theorem}

\begin{proof}
From the fact that $K_1(x) = Ax$ and $K_0(x) = A$ and the definition
of tensor product of complexes, we can identify $C_p(x)$ with $C_p \oplus C_{p-1}$, and
for $\xi \in C_p$, $\eta \in C_{p-1}$ we have
\[
  d(\xi, \eta) = (d\xi + (-1)^{p-1}x\eta, d\eta).
\]
The first assertion is clear from this. Moreover, $H_p(C'_\bullet) = H_{p-1}(C_\bullet)$ is also
clear, and if $\eta \in C_{p-1} = C'_p$ satisfies $d\eta = 0$ then in $C_\bullet(x)$ we have $d(0, \eta) =
((-1)^{p-1}x\eta, 0)$, so that the long exact sequence has the form indicated
in the theorem. Finally, if $d(\xi, \eta) = 0$ then $d\eta = 0$ and $d\xi = (-1)^px\eta$, so
that $x \cdot (\xi, \eta) = (x\xi, (-1)^px\eta) \in d(C_{p+1}(x))$, and therefore $x \cdot H_p(C_\bullet(x)) = 0$.
\end{proof}

Applying this theorem to $K_\bullet(x, M)$ and using the commutativity of
tensor product of complexes, we see that the ideal $(\mathfrak{x}) = (x_1, \ldots, x_n)$
generated by $\mathfrak{x}$ annihilates the homology groups $H_p(x, M)$:
\[
  (\mathfrak{x}) \cdot H_p(x, M) = 0 \quad \text{for all } p.
\]

\begin{theorem}\label{thm:16.5}
\begin{enumerate}[(i)]
  \item Let $A$ be a ring, $M$ an $A$-module, and $x_1, \ldots, x_n$ an $M$-sequence; then
        \[
          H_p(x, M) = 0 \text{ for } p > 0 \quad \text{and} \quad H_0(x, M) = M/\mathfrak{x}M.
        \]
  \item Suppose that one of the following two conditions ($\alpha$) or ($\beta$) holds:
        \begin{enumerate}[($\alpha$)]
          \item $(A, m)$ is a local ring, $x_1, \ldots, x_n \in m$ and $M$ is a finite $A$-module;
          \item $A$ is an $\N$-graded ring, $M$ is an $\N$-graded $A$-module, and $x_1, \ldots, x_n$ are
                homogeneous elements of degree $> 0$.
        \end{enumerate}
        Then the converse of (i) holds in the following strong form: if $H_p(x, M) = 0$
        and $M \neq 0$ then $x_1, \ldots, x_n$ is an $M$-sequence.
\end{enumerate}
\end{theorem}

\begin{proof}
We use induction on $n$.

(i) When $n = 1$ we have $H_1(x, M) = \{\xi \in M \mid x\xi = 0\} = 0$, so that
the assertion holds. When $n > 1$, for $p > 1$ the previous theorem provides
an exact sequence
\[
  0 = H_p(x_1, \ldots, x_{n-1}, M) \to H_p(x_1, \ldots, x_n, M)
  \to H_{p-1}(x_1, \ldots, x_{n-1}, M) = 0.
\]
so that $H_p(x_1, \ldots, x_n, M) = 0$. For $p = 1$, setting $M_i = M/(x_1, \ldots, x_i)M$ we
have an exact sequence
\[
  0 \to H_1(x_n, M_{n-1}) \to H_1(x_1, \ldots, x_n, M) \to H_1(x_1, \ldots, x_{n-1}, M) = M_{n-2}
\]
\[
  \xrightarrow{x_n} M_{n-1},
\]
and since $x_n$ is $M_{n-1}$-regular we have $H_1(x_n, M) = 0$.

(ii) Assuming either ($\alpha$) or ($\beta$), $M \neq 0$ implies that $M_i \neq 0$ for $1 \le i \le n$. By
hypothesis and by the previous theorem,
\[
  H_1(x_1, \ldots, x_{n-1}, M) \xrightarrow{x_n} H_1(x_1, \ldots, x_{n-1}, M) \to H_1(x, M) = 0;
\]
but quite generally $H_1(x_1, \ldots, x_{n-1}, M)$ is a finite $A$-module in case ($\alpha$), or an $\N$-graded
module in case ($\beta$), so that by NAK, $H_1(x_1, \ldots, x_{n-1}, M) = 0$. Thus by
induction $x_1, \ldots, x_{n-1}$ is an $M$-sequence. Now by the same exact sequence
and the first part of (i), we see that $x_n$ is $M_{n-1}$-regular, and therefore
$x_1, \ldots, x_n$ is an $M$-sequence.
\end{proof}

Let $A$ be a ring, $M$ an $A$-module and $I$ an ideal of $A$. If $a_1, \ldots, a_n$ are
elements of $I$, we say that they form a \emph{maximal $M$-sequence in $I$}\index{maximal M-sequence} if $a_1, \ldots, a_n$
is an $M$-sequence, and $a_1, \ldots, a_n, b$ is not an $M$-sequence for any $b \in I$. If
$a_1, \ldots, a_n$ is an $M$-sequence then $a_1M, (a_1, a_2)M, \ldots, (a_1, \ldots, a_n)M$ is strictly
increasing, so that the chain of ideals $(a_1) \subset (a_1, a_2) \subset \cdots$ is also strict-
ly increasing. If $A$ is Noetherian this cannot continue indefinitely,
so that any $M$-sequence can be extended until we arrive at a maximal
$M$-sequence.

In Theorems 6--8 below, the hypothesis that $M$ is a finite
$A$-module can be weakened to the statement that $M$ is a finite $B$-module for a
homomorphism $A \to B$ of Noetherian rings, as one sees on inspecting
the proof. The reason for this is that, if we set $\Ass_B(M) = \{P_1, \ldots, P_r\}$ and
$P_i \cap A = p_i$, then any ideal of $A$ consisting entirely of zero-divisors of $M$ is
contained in $\bigcup p_i$, and therefore contained in one of the $p_i$. Note that
according to [M], (9.A), we have $\Ass_A(M) = \{p_1, \ldots, p_r\}$.

\begin{theorem}\label{thm:16.6}
Let $A$ be a Noetherian ring, $M$ a finite $A$-module and $I$
an ideal of $A$; suppose that $IM \neq M$. For a given integer $n > 0$ the following
conditions are equivalent;
\begin{enumerate}[(1)]
  \item $\Ext^i_A(N, M) = 0$ for all $i < n$ and for any finite $A$-module $N$ with
        $\Supp(N) \subset V(I)$;
  \item $\Ext^i_A(A/I, M) = 0$ for all $i < n$;
  \item $\Ext^i_A(A/I, M) = 0$ for all $i < n$ and for some finite $A$-module $N$ with
        $\Supp(N) = V(I)$;
  \item there exists an $M$-sequence of length $n$ contained in $I$.
\end{enumerate}
\end{theorem}

\begin{proof}
$(1) \Rightarrow (2) \Rightarrow (3)$ are obvious. For $(2') \Rightarrow (3)$, if $I$ consists only of zero-
divisors of $M$ then there exists an associated prime $P$ of $M$ containing $I$
(this is where we need the finiteness of $M$). Hence there is an injective
map $A/P \to M$. Localising at $P$, we see that $\Hom_{A_P}(A_P/PA_P, M_P) \neq 0$, where
$(A/P)_P = A_P/PA_P$. Now $P \in V(I) = \Supp(N)$, so that $N_P \neq 0$, and hence by
NAK, $N_P/PN_P = N \otimes_A k \neq 0$. Thus $N \otimes k$ is a non-zero vector space
over $k$, and $\Hom_A(N \otimes k, k) \neq 0$. Putting together what we have said,
we can follow the composite $N_P \to N \otimes k \to k \to M_P$ to show that
$\Hom_{A_P}(N_P, M_P) \neq 0$. The left-hand side is equal to $(\Hom_A(N, M))_P$, so that
$\Hom_A(N, M) \neq 0$. But this contradicts $(2')$. Hence $I$ contains an
$M$-regular element $f$. By assumption, $M/IM \neq 0$, and if $n = 1$ then we are
done. If $n > 1$ we set $M_1 = M/fM$; then from the exact sequence
\[
  0 \to M \xrightarrow{f} M \to M_1 \to 0
\]
we get $\Ext^i_A(N, M_1) = 0$ for $i < n - 1$, so that by induction $I$ contains an
$M_1$-sequence $f_2, \ldots, f_n$.

For the proof of $(3) \Rightarrow (1)$ we do not need to assume that $A$ is Noetherian
or $M$ finite. Let $f_1, \ldots, f_n \in I$ be an $M$-sequence; we have the exact sequence
\[
  0 \to M \xrightarrow{f_1} M \to M_1 \to 0,
\]
and if $n > 1$ the inductive hypothesis $\Ext^i_A(N, M_1) = 0$ for $i < n - 1$, so that
\[
  0 \to \Ext^i_A(N, M) \xrightarrow{f_1} \Ext^i_A(N, M)
\]
is exact for $i < n$. But $\Ext^i_A(N, M)$ is annihilated by elements of $\ann(N)$.
Since $\Supp(N) = V(\ann(N)) \subset V(I)$, we have $I \subset \sqrt{\ann(N)}$, and a
sufficiently large power of $f_1$ annihilates $\Ext^i_A(N, M)$. Therefore,
$\Ext^i_A(N, M) = 0$ for $i < n$.
\end{proof}

Let $M$ and $I$ be as in the above theorem, and $a_1, \ldots, a_n$ an $M$-sequence
in $I$. For $1 \le i \le n$, set $M_i = M/(a_1, \ldots, a_i)M$; then it is easy to see
that $\Hom_A(A/I, M_i) \supset \Ext^1_A(A/I, M_{i-1}) \supset \cdots \supset \Ext^i_A(A/I, M)$. Therefore,
if $\Ext^i_A(A/I, M) = 0$ we can find another element $a_{n+1} \in I$ such that
$a_1, \ldots, a_{n+1}$ is an $M$-sequence. Hence if $a_1, \ldots, a_n$ is a maximal $M$-sequence
in $I$ we must have $\Ext^n_A(A/I, M) \neq 0$. We thus obtain the following theorem.

\begin{theorem}\label{thm:16.7}
Let $A$ be a Noetherian ring, $I$ an ideal of $A$ and $M$ a finite
$A$-module such that $M \neq IM$; then the length of a maximal $M$-sequence
in $I$ is a well-determined integer $n$, and $n$ is determined by
\[
  \Ext^i_A(A/I, M) = 0 \text{ for } i < n \quad \text{and} \quad \Ext^n_A(A/I, M) \neq 0.
\]
We write $n = \depth(I, M)$, and call $n$ the \emph{$I$-depth}\index{depth} of $M$. (If $M = IM$, the
$I$-depth is by convention $\infty$.) Theorem 7 takes the form
\[
  \depth(I, M) = \inf\{i \mid \Ext^i_A(A/I, M) \neq 0\}.
\]
In particular for a Noetherian local ring $(A, m, k)$, we call $\depth(m, M)$
simply the \emph{depth} of $M$, and write $\depth M$ or $\depth_A M$:
\[
  \depth M = \inf\{i \mid \Ext^i_A(k, M) \neq 0\}.
\]
From Theorem 6 we see that if $V(I) = V(I')$ then $\depth(I, M) =
\depth (I', M)$; this also follows easily from Theorem 1.

If $\ann (M) = \mathfrak{a}$ and we set $\bar{A} = A/\mathfrak{a}$ then $M$ is also an $\bar{A}$-module. Writing
$\bar{a}$ or $\bar{I}$ for the image of an element $a$ or an ideal $I$ of $A$ under the natural
homomorphism $A \to \bar{A}$ we clearly have that $a_1, \ldots, a_n$ is an $M$-sequence
if and only if $\bar{a}_1, \ldots, \bar{a}_n$ is. Thus $\depth (I, M) = \depth (\bar{I}, M)$, and if we set
$I + \mathfrak{a} = J$, then since $\bar{I} = \bar{J}$ we also have $\depth (I, M) = \depth (J, M)$.
\end{theorem}

We can also prove that the length of a maximal $M$-sequence is
well-determined by means of the Koszul complex.

\begin{theorem}\label{thm:16.8}
Let $A$ be a Noetherian ring, $I = (y_1, \ldots, y_n)$ an ideal of $A$,
and $M$ a finite $A$-module such that $M \neq IM$. If we set
\[
  q = \sup\{i \mid H_i(y, M) \neq 0\},
\]
then any maximal $M$-sequence in $I$ has length $n - q$.
\end{theorem}

\begin{proof}
Let $x_1, \ldots, x_s$ be a maximal $M$-sequence in $I$; we argue by induction
on $s$. If $s = 0$ then every element of $I$ is a zero-divisor of $M$, so that there
exists $P \in \Ass(M)$ containing $I$. By definition of $\Ass$, there exists $0 \neq \xi \in M$
such that $P = \ann(\xi)$, and hence $I\xi = 0$. Thus $\xi \in H_n(y, M)$ so that $q = n$,
and the assertion holds in this case.

If $s > 0$ we set $M_1 = M/x_1M$; then from the exact sequence
\[
  0 \to M \xrightarrow{x_1} M \to M_1 \to 0
\]
and from the fact that $IH_i(y, M) = 0$ (by Theorem 4), it follows that
\[
  0 \to H_i(y, M) \to H_i(y, M_1) \to H_{i-1}(y, M) \to 0
\]
is exact for every $i$. Thus $H_{q+1}(y, M_1) \neq 0$ and $H_i(y, M_1) = 0$ for $i > q + 1$;
but $x_2, \ldots, x_s$ is a maximal $M_1$-sequence in $I$, so that by induction we
have $q + 1 = n - (s - 1)$, and therefore $q = n - s$.
\end{proof}

In other words, $\depth(I, M)$ is the number of successive zero terms from
the left in the sequence
\[
  H_n(y, M), H_{n-1}(y, M), \ldots, H_0(y, M) = M/IM \neq 0.
\]
This fact is sometimes referred to as the `depth sensitivity' of the Koszul
complex.

\begin{corollary*}
In the situation of the theorem, $y_1, \ldots, y_n$ is an $M$-sequence if
and only if $\depth (I, M) = n$.
\end{corollary*}

\begin{proof}
$\depth(I, M) = n \iff H_i(y, M) = 0$ for all $i > 0 \iff y$ is an $M$-sequence.
\end{proof}

\subsection*{Grade}

A little before Auslander and Buchsbaum [2], Rees [5] introduced and
developed the theory of another notion related to regular sequences, that of
grade.\index{grade} Let $A$ be a Noetherian ring and $M \neq 0$ a finite $A$-module. Then Rees
made the definition
\[
  \grade M = \inf\{i \mid \Ext^i_A(M, A) \neq 0\}.
\]
For a proper ideal $J$ of $A$ we also call $\grade (A/J)$ the \emph{grade of the ideal}
$J$, and write $\grade J$. If we set $\mathfrak{a} = \ann(M)$ then since $\Supp(M) = V(\mathfrak{a})$,
Theorem 6 gives $\grade M = \depth(\mathfrak{a}, A)$. Moreover, if $g = \grade M$ then
$\Ext^g_A(M, A) \neq 0$, so that
\[
  \grade M \le \projdim M.
\]
If $I$ is an ideal then $\grade I (=\grade(A/I)) = \depth(I, A)$ is the length
of a maximal $A$-sequence in $I$, but in general if $a_1, \ldots, a_r$ is an $A$-sequence
then one sees easily from Theorem~\ref{thm:13.5} that $\htop(a_1, \ldots, a_r) = r$. Thus if
$a_1, \ldots, a_r$ is a maximal $A$-sequence in $I$, we have $r = \htop(a_1, \ldots, a_r) \le \htop I$.
Hence for an ideal $I$ we have $\grade I \le \htop I$.

\begin{theorem}\label{thm:16.9}
Let $A$ be a Noetherian ring, and $M, N$ finite $A$-modules;
suppose $M \neq 0$, $\grade M = k$ and $\projdim N = l \le k$. Then
\[
  \Ext^i_A(M, N) = 0 \quad \text{for } i < k - l.
\]
\end{theorem}

\begin{proof}
We use induction on $l$. If $l = 0$ then $N$ is a direct summand of some
free module $A^m$, so that we need only say what happens for $N = A$, but
then the assertion is just the definition of grade. If $l > 0$ we choose an
exact sequence
\[
  0 \to N_1 \to L_0 \to N \to 0
\]
with $L_0$ a finite free module; then $\projdim N_1 = l - 1$, so that by induction
\[
  \Ext^i_A(M, L_0) = 0 \quad \text{for } i < k \quad \text{and}
\]
\[
  \Ext^i_A(M, N_1) = 0 \quad \text{for } i < k - l;
\]
the assertion follows from this.
\end{proof}

% ============ 习题 ============
\section*{Exercises to \S\ref{sec:16}}
\emph{Prove the following propositions.}

\begin{exercise}\label{ex:16.1-ex}
Let $(A, m)$ be a Noetherian local ring, $M \neq 0$ a finite $A$-module, and
$a_1, \ldots, a_r \in m$ an $M$-sequence. Set $M' = M/(a_1, \ldots, a_r)M$. Then
\[
  \dim M' = \dim M - r.
\]
\end{exercise}

\begin{exercise}\label{ex:16.2-ex}
Let $A$ be a Noetherian ring, $\mathfrak{a}$ and $\mathfrak{b}$ ideals of $A$; then if $\grade \mathfrak{a}
> \projdim A/\mathfrak{b}$ we have $\mathfrak{b}:\mathfrak{a} = \mathfrak{b}$.
\end{exercise}

\begin{exercise}\label{ex:16.3-ex}
Let $A$ be a Noetherian ring. A proper ideal $I$ of $A$ is called a \emph{perfect ideal} if
$\grade I = \projdim A/I$. If $I$ is a perfect ideal of grade $k$ then all the prime
divisors of $I$ have grade $k$.
\end{exercise}

\begin{remark}
Quite generally, we have $\grade I (=\grade(A/I)) \le \projdim A/I$. If $A$ is a
regular local ring and $P \in \Spec A$ then as we will see in Theorems 19.1 and 19.2, $P$ is
perfect, that is $A/P$ is Cohen--Macaulay.
\end{remark}

\begin{exercise}\label{ex:16.4-ex}
Let $f\colon A \to B$ be a flat ring homomorphism, $M$ an $A$-module, and
$a_1, \ldots, a_n \in A$ an $M$-sequence; if $(M/(a_1, \ldots, a_n)M) \otimes_A B \neq 0$ then
$f(a_1), \ldots, f(a_n)$ is an $M \otimes_A B$-sequence.
\end{exercise}

\begin{exercise}\label{ex:16.5-ex}
Let $A$ be a Noetherian local ring, $M$ a finite $A$-module, and $P$ a prime ideal
of $A$; show that $\depth(P, M) \le \depth_{A_P} M_P$, and construct an example
where the inequality is strict.
\end{exercise}

\begin{exercise}\label{ex:16.6-ex}
Let $A$ be a ring and $a_1, \ldots, a_n \in A$ an $A$-quasi-regular sequence. If $A$
contains a field $k$ then $a_1, \ldots, a_n$ are algebraically independent over $k$.
\end{exercise}

\begin{exercise}\label{ex:16.7-ex}
Let $(A, m)$ and $(B, n)$ be Noetherian local rings, and suppose that $A \subset B$,
$n \cap A = m$ and that $mB$ is an $n$-primary ideal. Then for a finite $B$-module
$M$ we have
\[
  \depth_A M = \depth_B M.
\]
\end{exercise}

\begin{exercise}\label{ex:16.8-ex}
Let $A$ be a ring, $P_1, \ldots, P_r$ prime ideals, $I$ an ideal, and $x$ an element of $A$. If
$x \notin I + \bigcup_{i=1}^{r} P_i$ then there is a $y \in I$ such that $x + y \notin \bigcup_{i=1}^{r} P_i$
(E. Davis).
\end{exercise}

\begin{exercise}\label{ex:16.9-ex}
Use the previous question to show the following: let $A$ be a Noetherian
ring, and suppose that $I \neq A$ is an ideal generated by $n$ elements; then
$\grade I \le n$, and if $\grade I = n$ then $I$ can be generated by an $A$-sequence
([K], Th. 125).
\end{exercise}

\begin{exercise}\label{ex:16.10-ex}
Let $A$ be a Noetherian ring, and suppose that $P$ is a height $r > 0$ prime
ideal generated by $r$ elements $a_1, \ldots, a_r$.
\begin{enumerate}[(i)]
  \item Suppose either that $A$ is local, or that $A$ is $\N$-graded and the $a_i$ are
        homogeneous of positive degree. Then $A$ is an integral domain, and for
        $1 \le i \le r$ the ideal $(a_1, \ldots, a_i)$ is prime; hence $a_1, \ldots, a_r$ is an $A$-sequence.
  \item In general $a_1, \ldots, a_r$ does not have to be an $A$-sequence, but $P$ can in
        any case be generated by an $A$-sequence (E. Davis).
\end{enumerate}
\end{exercise}

% ============ §17 Cohen-Macaulay rings (原书 pp.133-138) ============
\section{Cohen--Macaulay rings}\label{sec:17}\index{Cohen-Macaulay rings, section 17}

\begin{theorem}[Ischebeck]\label{thm:17.1}
Let $(A, m)$ be a Noetherian local ring, $M$ and $N$
non-zero finite $A$-modules, and suppose that $\depth M = k$, $\dim N = r$. Then
\[
  \Ext^i_A(N, M) = 0 \quad \text{for } i < k - r.
\]
\end{theorem}

\begin{proof}
By induction on $r$; if $r = 0$ then $\Supp(N) = \{m\}$ and the assertion
holds by Theorem~\ref{thm:16.6}. Suppose $r > 0$. By Theorem~\ref{thm:6.4}, there exists a
chain
\[
  N = N_0 \supset N_1 \supset \cdots \supset N_t = (0) \quad \text{with } N_j/N_{j+1} \cong A/P_j
\]
of submodules $N_j$, where $P_j \in \Spec A$. It is easy to see that if $\Ext^i_A
(N_j/N_{j+1}, M) = 0$ for each $j$ then $\Ext^i_A(N, M) = 0$, and since $\dim N_j/N_{j+1} \le
\dim N = r$ it is enough to prove that $\Ext^i_A(N, M) = 0$ for $i < k - r$ in the
case $N = A/P$ with $P \in \Spec A$ and $\dim N = r$. Since $r > 0$ we can take
an element $x \in m - P$ and get the exact sequence
\[
  0 \to N \xrightarrow{x} N \to N' \to 0,
\]
where $N' = A/(P, x)$; then $\dim N' < r$ so that by induction we have
$\Ext^i_A(N', M) = 0$ for $i < k - r + 1$. Thus for $i < k - r$ we have an exact
sequence
\[
  0 \to \Ext^i_A(N, M) \xrightarrow{x} \Ext^i_A(N, M) \to \Ext^{i+1}_A(N', M) = 0.
\]
We have $x \in m$ so that by NAK, $\Ext^i_A(N, M) = 0$.
\end{proof}

\begin{theorem}\label{thm:17.2}
Let $A$ be a Noetherian local ring, $M$ a finite $A$-module,
and assume that $P \in \Ass(M)$; then $\dim(A/P) \ge \depth M$.
\end{theorem}

\begin{proof}
If $P \in \Ass(M)$ then $\Hom_A(A/P, M) \neq 0$, so that by the previous
theorem we cannot have $\dim A/P < \depth M$.
\end{proof}

\begin{definition}
Let $(A, m, k)$ be a Noetherian local ring, and $M$ a finite
$A$-module. We say that $M$ is a \emph{Cohen--Macaulay module}\index{Cohen-Macaulay module} (abbreviated to
CM module) if $M \neq 0$ and $\depth M = \dim M$, or if $M = 0$. If $A$ itself is a
CM module we say that $A$ is a \emph{CM ring}\index{Cohen-Macaulay ring} or a \emph{Macaulay ring}.
\end{definition}

\begin{theorem}\label{thm:17.3}
Let $A$ be a Noetherian local ring and $M$ a finite $A$-module.
\begin{enumerate}[(i)]
  \item If $M$ is a CM module then for any $P \in \Ass(M)$ we have
        $\dim(A/P) = \dim M = \depth M$. Hence $M$ has no embedded associated
        primes.
  \item If $a_1, \ldots, a_r \in m$ is an $M$-sequence and we set $M' = M/(a_1, \ldots, a_r)M$ then
        \[
          M \text{ is a CM module} \iff M' \text{ is a CM module}.
        \]
  \item If $M$ is a CM module then $M_P$ is a CM module over $A_P$ for
        every $P \in \Spec A$, and if $M_P \neq 0$ then
        \[
          \depth (P, M) = \depth_{A_P} M_P.
        \]
\end{enumerate}
\end{theorem}

\begin{proof}
(i) Quite generally, we have
\[
  \dim M = \sup\{\dim A/P \mid P \in \Ass M\}
\]
\[
  \ge \inf\{\dim A/P \mid P \in \Ass M\} \ge \depth M,
\]
so that this is clear.

(ii) By definition $\depth M' = \depth M - r$, and by Ex.~\ref{ex:16.1-ex}, $\dim M' =
\dim M - r$, so that this is clear.

(iii) It is enough to consider the case $M_P \neq 0$, when $P \supset \ann(M)$. Then
quite generally we have $\dim M_P \ge \depth M_P \ge \depth(P, M)$, so that we need
only show that
\[
  \dim M_P = \depth (P, M).
\]
We prove this by induction on $\depth(P, M)$. If $\depth(P, M) = 0$ then $P$ is
contained in an associated prime of $M$, but in view of $P \supset \ann(M)$ and
the fact that by (i) all the associated primes of $M$ are minimal, it follows
that $P$ is itself an associated prime of $M$; therefore $\dim M_P = 0$. If
$\depth(P, M) > 0$ then we can take an $M$-regular element $a \in P$, and set
$M' = M/aM$. Then
\[
  \depth (P, M') = \depth (P, M) - 1,
\]
$M'$ is a CM module with $M'_P \neq 0$, so that by induction $\dim M'_P =
\depth(P, M) - 1$; moreover, $a$ is $M_P$-regular as an element of $A_P$, and
so that using Ex.~\ref{ex:16.1-ex} once more, we have $\dim M_P =
\dim M'_P + 1$. Putting these together gives $\depth (P, M) = \dim M_P$.
\end{proof}

\begin{theorem}\label{thm:17.4}
Let $(A, m)$ be a CM local ring.
\begin{enumerate}[(i)]
  \item For a proper ideal $I$ of $A$ we have
        \[
          \htop I = \depth (I, A) = \grade I, \quad \text{and} \quad \htop I + \dim A/I = \dim A.
        \]
  \item For elements $a_1, \ldots, a_r \in m$ the following four conditions are
        equivalent:
        \begin{enumerate}[(1)]
          \item $a_1, \ldots, a_r$ is an $A$-sequence;
          \item $\htop(a_1, \ldots, a_i) = i$ for $1 \le i \le r$;
          \item $\htop(a_1, \ldots, a_r) = r$;
          \item $a_1, \ldots, a_r$ is part of a system of parameters of $A$.
        \end{enumerate}
  \item If $P \subset Q$ are prime ideals of $A$ then
        \[
          \htop Q - \htop P = \htop (Q/P).
        \]
\end{enumerate}
\end{theorem}

\begin{proof}
(iii) The implication $(1) \Rightarrow (2)$ follows from Theorem~\ref{thm:13.5}, together with
the fact that from the definition of $A$-sequence we have $0 < \htop(a_1) \le
\htop(a_1, a_2) \le \cdots \le r$.

$(2) \Rightarrow (3)$ is obvious.

$(3) \Rightarrow (4)$ If $\dim A = r$ this is obvious; if $\dim A > r$ then $m$ is not a minimal
prime divisor of $(a_1, \ldots, a_r)$, so that we can choose $a_{r+1} \in m$ not contained
in any minimal prime divisor of $(a_1, \ldots, a_r)$, and then $\htop(a_1, \ldots, a_{r+1}) = r + 1$.
Proceeding in the same way we arrive at a system of parameters of $A$.
(Note that so far we have not used the CM assumption.)

$(4) \Rightarrow (1)$ It is enough to show that any system of parameters $x_1, \ldots, x_n$
($n = \dim A$) is an $A$-sequence. If $P \in \Ass(A)$ then by Theorem 3, (i),
$\dim A/P = n$, so that $x_1 \notin P$. Thus $x_1$ is $A$-regular. Therefore if we set
$A' = A/x_1A$ we have by the previous theorem that $A'$ is an $(n - 1)$-
dimensional CM ring, and the images of $x_2, \ldots, x_n$ form a system of
parameters of $A'$. Thus by induction on $n$ we see that $x_1, \ldots, x_n$ is an
$A$-sequence.

(i) If $\htop I = r$ then we can take $a_1, \ldots, a_r \in I$ such that $\htop(a_1, \ldots, a_i) = i$ for
$1 \le i \le r$. Thus by (iii), $a_1, \ldots, a_r$ is an $A$-sequence. Thus $r \le \grade I$.
Conversely, if $b_1, \ldots, b_s$ is an $A$-sequence in $I$ then $\htop(b_1, \ldots, b_s) = s \le \htop I$, and
hence $r \ge \grade I$, so that equality must hold. For the second equality, let
$S$ be the set of minimal prime divisors of $I$, we have
\[
  \htop I = \inf\{\htop P \mid P \in S\}
\]
and
\[
  \dim(A/I) = \sup\{\dim A/P \mid P \in S\},
\]
so it is enough to show that $\htop P = \dim A - \dim A/P$ for every $P \in S$.
Let $\htop P = r$ and $\dim A = n$. By Theorem 3, (iii), $A_P$ is a CM ring
and $r = \depth(P, A)$. Now if we take an $A$-sequence $a_1, \ldots, a_r \in P$ then
by Theorem 3, (ii), $A/(a_1, \ldots, a_r)$ is an $(n - r)$-dimensional CM ring, and from
the fact that $\htop (a_1, \ldots, a_r) = r = \htop P$ we see that $P$ is a minimal prime
divisor of $(a_1, \ldots, a_r)$; thus by Theorem 3, (i), $\dim A/P = \dim A/(a_1, \ldots, a_r) =
n - r$.

(ii) Let $P \subset Q$ be prime ideals of $A$. Then since $A_Q$ is a CM ring, (i)
above gives $\dim A_Q = \htop QA_Q + \dim A_Q/QA_Q$; in other words $\htop P - \htop Q =
\htop (Q/P)$.
\end{proof}

If one system of parameters of a Noetherian local ring $A$ is an $A$-sequence
then $\depth A = \dim A$, so that $A$ is a CM ring, and therefore, by the above
theorem, every system of parameters of $A$ is an $A$-sequence.

\begin{theorem}\label{thm:17.5}
Let $A$ be a Noetherian local ring and $\hat{A}$ its completion;
then
\begin{enumerate}[(i)]
  \item $\depth A = \depth \hat{A}$;
  \item $A$ is CM $\iff \hat{A}$ is CM.
\end{enumerate}
\end{theorem}

\begin{proof}
(i) This comes for example from the fact that
\[
  \Ext^i_A(A/m, A) \otimes_A \hat{A} = \Ext^i_{\hat{A}}(\hat{A}/m\hat{A}, \hat{A}) \quad \text{for all } i.
\]
(ii) follows from (i) and the fact that
$\dim A = \dim \hat{A}$.
\end{proof}

\begin{definition}
A proper ideal $I$ in a Noetherian ring $A$ is said to be \emph{unmixed}\index{unmixed ideal}
if the heights of its prime divisors are all equal. We say that the \emph{unmixedness}
\emph{theorem} holds for $A$ if for every $r > 0$, every height $r$ ideal $I$ of $A$ generated
by $r$ elements is unmixed. This includes as the case $r = 0$ the statement
that $(0)$ is unmixed. By Theorem~\ref{thm:13.5}, if $I$ is an ideal satisfying the
hypotheses of this proposition, then all the minimal prime divisors of $I$
have height $r$, so that to say that $I$ is unmixed is to say that $I$ does not
have embedded prime divisors.
\end{definition}

A Noetherian ring $A$ is said to be a \emph{CM ring}\index{Cohen-Macaulay ring} if $A_m$ is a CM local
ring for every maximal ideal $m$ of $A$. By Theorem 3, (iii), a localisation
$S^{-1}A$ of a CM ring $A$ is again CM.

\begin{theorem}\label{thm:17.6}
A necessary and sufficient condition for a Noetherian ring
$A$ to be a CM ring is that the unmixedness theorem holds for $A$.
\end{theorem}

\begin{proof}
First suppose that $A$ is a CM ring and that $I = (a_1, \ldots, a_r)$ is an
ideal of $A$ with $\htop I = r$. We assume that $P$ is an embedded prime divisor
of $I$ and derive a contradiction. Localising at $P$ we can assume that $A$ is
a CM local ring; then by Theorem 4, (iii), $a_1, \ldots, a_r$ is an $A$-sequence, and
hence $A/I$ is also a CM local ring. But then $I$ does not have embedded
prime divisors, and this is a contradiction. Next we suppose that the
unmixedness theorem holds for $A$. If $P \in \Spec A$ with $\htop P = r$ then we can
choose $a_1, \ldots, a_r \in P$ such that
\[
  \htop(a_1, \ldots, a_i) = i \quad \text{for } 1 \le i \le r.
\]
Then by the unmixedness theorem, all the prime divisors of $(a_1, \ldots, a_i)$
have height $i$, and therefore do not contain $a_{i+1}$. Hence $a_{i+1}$ is an
$A/(a_1, \ldots, a_i)$-regular element; in other words, $a_1, \ldots, a_r$ is an $A$-sequence.
Therefore $\depth A_P = r = \dim A_P$, so that $A_P$ is a CM local ring; $P$ was
any element of $\Spec A$, so that $A$ is a CM ring.
\end{proof}

The unmixedness theorem for polynomial rings over a field was a
brilliant early result of Macaulay in 1916; for regular local rings, the
unmixedness theorem was proved by I. S. Cohen [1] in 1946. This explains
the term Cohen--Macaulay. Having come this far, we are now in a position
to give easy proofs of these two theorems.

\begin{theorem}\label{thm:17.7}
If $A$ is a CM ring then so is $A[X_1, \ldots, X_n]$.
\end{theorem}

\begin{proof}
We need only consider the case $n = 1$. Set $B = A[X]$ and let $P$ be
a maximal ideal of $B$. Set $P \cap A = m$; then $B_P$ is also a localisation of
$A_m[X]$, so that replacing $A$ by $A_m$ we have a local CM ring
$A$ with maximal ideal $m$, and we need only prove that $B_P$ is CM. Setting
$A/m = k$ we get
\[
  B/mB = k[X],
\]
so that $P/mB$ is a principal ideal of $k[X]$ generated by an irreducible
monic polynomial $\bar{\varphi}(X)$. If we let $f(X) \in A[X]$ be a monic polynomial of
$A[X]$ which reduces to $\varphi(X)$ modulo $mB$ then $P = (m, f)$. We choose a
system of parameters $a_1, \ldots, a_n$ for $A$, so that $a_1, \ldots, a_n, f$ is a system of
parameters of $B_P$. Since $B$ is flat over $A$ the $A$-sequence $a_1, \ldots, a_n$ is also
a $B$-sequence. We set $A/(a_1, \ldots, a_n) = A'$; then the image of $f$ in $A'[X]$ is
a monic polynomial, and therefore $A'[X]$-regular, so that $a_1, \ldots, a_n, f$ is
a $B$-sequence, and
\[
  \depth B_P \ge \depth (P, B) \ge n + 1 = \dim B_P.
\]
Therefore $B_P$ is a CM ring.
\end{proof}

\begin{remark}
If $A$ is a CM local ring, then a similar (if anything, rather easier)
method can be used to prove that $A[[X]]$ is also CM. The statement also
holds for a non-local CM ring, but the proof is a little more complicated,
and we leave it to \S23.
\end{remark}

\begin{theorem}\label{thm:17.8}
A regular local ring is a CM ring.
\end{theorem}

\begin{proof}
Let $(A, m)$ be an $n$-dimensional regular local ring, and $x_1, \ldots, x_n$ a
regular system of parameters. By Theorems 14.2 and 14.3, $(x_1), (x_1, x_2), \ldots,
(x_1, \ldots, x_n)$ is a strictly increasing chain of prime ideals; therefore $x_1, \ldots, x_n$
is an $A$-sequence.
\end{proof}

\begin{theorem}\label{thm:17.9}
Any quotient of a CM ring is universally catenary.
\end{theorem}

\begin{proof}
Clear from Theorems 7 and 4.
\end{proof}

\begin{theorem}\label{thm:17.10}
A necessary and sufficient condition for a Noetherian local
ring $(A, m, k)$ to be a regular ring is that $\gr_m(A)$ is isomorphic as a graded $k$-
algebra to a polynomial ring over $k$.
\end{theorem}

\begin{proof}
If $A$ is regular, let $x_1, \ldots, x_r$ be a regular system of parameters, that
is a minimal basis of $m$; then $x_1, \ldots, x_r$ is an $A$-sequence, so that by
Theorem~\ref{thm:16.2} (see also Theorem~\ref{thm:14.4} for another proof) $\gr_m(A) \cong
k[X_1, \ldots, X_r]$. Conversely, if $\gr_m(A) \cong k[X_1, \ldots, X_r]$, then comparing
the homogeneous components of degree $1$, we see that $m/m^2 \cong kX_1
+ \cdots + kX_r$. On the other hand, the homogeneous component of degree $n$
of $k[X_1, \ldots, X_r]$ is a vector space over $k$ of dimension $\binom{n+r-1}{r-1}$,
so that the Samuel function is
\[
  \chi_A(n) = l(A/m^{n+1}) = \sum_{i=0}^{n} l(m^i/m^{i+1}) = \binom{n+r}{r},
\]
and $\dim A = r$. Therefore $A$ is regular.
\end{proof}

We can also characterise CM local rings in terms of properties of
multiplicities. Let $A$ be a Noetherian local ring. An ideal of $A$ is said to
be a \emph{parameter ideal}\index{parameter ideal} if it can be generated by a system of parameters. By
Theorem~\ref{thm:14.10}, if $q$ is a parameter ideal then $l(A/q) \ge e(q)$. As we are about
to see, equality here is characteristic of CM rings.

\begin{theorem}\label{thm:17.11}
The following three conditions on a Noetherian local ring
$(A, m)$ are equivalent:
\begin{enumerate}[(1)]
  \item $A$ is a CM ring;
  \item $l(A/q) = e(q)$ for any parameter ideal $q$ of $A$;
  \item $l(A/q) = e(q)$ for some parameter ideal $q$ of $A$.
\end{enumerate}
\end{theorem}

\begin{proof}
$(1) \Rightarrow (2)$. If $x_1, \ldots, x_d$ is a system of parameters of $A$ and $q =
(x_1, \ldots, x_d)$ then by Theorem~\ref{thm:16.2}, $\gr_q(A) \cong (A/q)[X_1, \ldots, X_d]$, so that
as in the proof of the previous theorem, $\chi_A(n) = l(A/q)\binom{n+d}{d}$, so that
$e(q) = l(A/q)d!$.

$(2) \Rightarrow (3)$ is obvious.

$(3) \Rightarrow (1)$ Suppose that $q = (x_1, \ldots, x_d)$ is a parameter ideal satisfying
$e(q) = l(A/q)$. We set $B = (A/q)[X_1, \ldots, X_d]$; then there is a homogeneous
ideal $\mathfrak{b}$ of $B$ such that $\gr_q(A) \cong B/\mathfrak{b}$. We write $\varphi_B(n)$ and $\varphi_{\mathfrak{b}}(n)$ for the Hilbert
polynomials of $B$ and $\mathfrak{b}$ (see \S\ref{sec:13}); then
\[
  l(q^n/q^{n+1}) = \varphi_B(n) - \varphi_{\mathfrak{b}}(n),
\]
and for $n \gg 0$ we have $l(q^n/q^{n+1}) = \chi_A(n) - \chi_A(n-1)$. The left-hand side is a
polynomial in $n$ of degree $d - 1$, and the coefficient of $n^{d-1}$ is $e(q)/(d - 1)!$. By
hypothesis $e(q) = l(A/q)$, so that $\varphi_{\mathfrak{b}}(n)$ must be a polynomial in $n$ of degree at
most $d - 2$. However, if $\mathfrak{b} \neq (0)$ then we can take a non-zero homogeneous
element $f(X) \in \mathfrak{b}$. If $m' \subset q$ and we set $\bar{m} = m/q$ then in $B$ we have
$\bar{m}^\nu = (0)$, and therefore replacing $f$ by the product of $f$ with a suitable
element of $\bar{m}$, we can assume that $f \neq 0$ but $\bar{m}f = 0$. Then
if $f$ has degree $p$,
\[
  \varphi_{\mathfrak{b}}(n) \ge \binom{n-p+d-1}{d-1},
\]
the length of the homogeneous component of degree $n - p$ in $(A/m)[X_1, \ldots, X_d]$. This
contradicts $\deg \varphi_{\mathfrak{b}} < d - 1$. Hence $\mathfrak{b} = (0)$, and
\[
  \gr_q(A) = B = (A/q)[X_1, \ldots, X_d],
\]
so that by Theorem~\ref{thm:16.3}, $\{x_1, \ldots, x_d\}$ is an $A$-sequence. Therefore $A$ is a CM
ring.
\end{proof}

% ============ 习题 ============
\section*{Exercises to \S\ref{sec:17}}
\emph{Prove the following propositions.}

\begin{exercise}\label{ex:17.1-ex}
\begin{enumerate}[(a)]
  \item A zero-dimensional Noetherian ring is a CM ring.
  \item A one-dimensional ring is CM provided that it is reduced ($=$ no
        nilpotent elements); also, construct an example of a one-dimensional ring
        which is not CM.
\end{enumerate}
\end{exercise}

\begin{exercise}\label{ex:17.2-ex}
Let $k$ be a field, $x, y$ indeterminates over $k$, and set
\[
  A = k[x^3, x^2y, xy^2, y^3] \subset k[x, y] \quad \text{and} \quad P = (x^3, x^2y, xy^2, y^3)A.
\]
Is $R = A_P$ a CM ring? How about
$k[x^4, x^3y, xy^3, y^4]$?
\end{exercise}

\begin{exercise}\label{ex:17.3-ex}
A two-dimensional normal ring is CM.
\end{exercise}

\begin{exercise}\label{ex:17.4-ex}
Let $A$ be a CM ring, $a_1, \ldots, a_r$ an $A$-sequence, and set $J = (a_1, \ldots, a_r)$. Then
for every integer $\nu$ the ring $A/J^\nu$ is CM, and therefore $J^\nu$ is unmixed.
\end{exercise}

\begin{exercise}\label{ex:17.5-ex}
Let $A$ be a Noetherian local ring and $P \in \Spec A$. Then
\begin{enumerate}[(i)]
  \item $\depth A \le \depth (P, A) + \dim A/P$;
  \item call $\dim A - \depth A$ the \emph{codepth} of $A$. Then $\codepth A \ge
        \codepth A_P$.
\end{enumerate}
\end{exercise}

\begin{exercise}\label{ex:17.6-ex}
Let $A$ be a Noetherian ring, $P \in \Spec A$ and set $G = \gr_P(A)$. If $G$ is an
integral domain then $P^n = P^{(n)}$ for all $n > 0$. (This observation is due to
Robbiano. One sees from it that if $P$ is a prime ideal generated by an $A$-
sequence then $P^n = P^{(n)}$.)
\end{exercise}

% ============ §18 Gorenstein rings (原书 pp.139-152) ============
\section{Gorenstein rings}\label{sec:18}\index{Gorenstein rings, section 18}

\begin{lemma}\label{lem:18.1}
Let $A$ be a ring, $M$ an $A$-module, and $n > 0$ a given integer,
\[
  \injdim M \le n \iff \Ext^{n+1}_A(A/I, M) = 0 \text{ for all ideals } I.
\]
If $A$ is Noetherian, then we can replace `for all ideals' by `for all prime
ideals' in the right-hand condition.
\end{lemma}

\begin{proof}
($\Rightarrow$) This is clear on calculating the Ext by an injective resolution
of $M$.

($\Leftarrow$) If $n = 0$ then from the exact sequence $0 \to I \to A \to A/I \to 0$ and
from the fact that $\Ext^1_A(A/I, M) = 0$ we get that $\Hom(A, M) \to
\Hom(I, M) \to 0$ is exact. Since this holds for every $I$, Theorem B3 of
Appendix B implies that $M$ is injective. Suppose then that $n > 0$.
There exists an exact sequence
\[
  0 \to M \to Q^0 \to Q^1 \to \cdots \to Q^{n-1} \to C \to 0,
\]
with each $Q^i$ injective. (We can obtain this by taking an injective resolution
of $M$ up to $Q^{n-1}$ and setting $C$ for the cokernel of $Q^{n-2} \to Q^{n-1}$.)
One sees easily that $\Ext^{n+1}_A(A/I, M) \cong \Ext^1_A(A/I, C)$, so that by the argument
used in the $n = 0$ case, $C$ is injective, and so $\injdim M \le n$.

If $A$ is Noetherian then by Theorem~\ref{thm:6.4}, any finite $A$-module $N$ has a
chain $N = N_0 \supset N_1 \supset \cdots \supset N_{r+1} = 0$ of submodules such that $N_j/N_{j+1} \cong
A/P_j$, with $P_j \in \Spec A$. Using this, if $\Ext^i_A(A/P, M) = 0$ for all prime ideals
$P$ then we also have $\Ext^i_A(N, M) = 0$ for all finite $A$-modules $N$. Now we
just have to apply this with $i = n + 1$ and $N = A/I$.
\end{proof}

\begin{lemma}\label{lem:18.2}
Let $A$ be a ring, $M$ and $N$ two $A$-modules, and $x \in A$;
suppose that $x$ is both $A$-regular and $M$-regular, and that $xN = 0$. Set $B = A/xA$ and
$\bar{M} = M/xM$. Then
\begin{enumerate}[(i)]
  \item $\Hom_B(N, \bar{M}) = 0$, and $\Ext^{n+1}_A(N, M) \cong \Ext^n_B(N, \bar{M})$ for all $n \ge 0$;
  \item $\Ext^n_A(M, N) \cong \Ext^n_B(\bar{M}, N)$ for all $n \ge 0$;
  \item $\Tor^A_n(M, N) \cong \Tor^B_n(\bar{M}, N)$ for all $n \ge 0$.
\end{enumerate}
\end{lemma}

\begin{proof}
(i) The first formula is obvious. For the second, set $T^n(N) =
\Ext^{n+1}_A(N, M)$, and view $T^n$ as a contravariant functor from the category
of $B$-modules to that of Abelian groups. Then first of all, the exact sequence
\[
  0 \to M \xrightarrow{x} M \to \bar{M} \to 0
\]
gives $T^0(N) = \Hom_A(N, M) = \Hom_B(N, \bar{M})$. Moreover, since $x$ is
$A$-regular we have $\projdim_A B = 1$, and therefore $T^n(B) = 0$ for $n > 0$,
so that $T^n(L) = 0$ for $n > 0$ and every projective $B$-module $L$. Finally, for
any short exact sequence $0 \to N' \to N \to N'' \to 0$ of $B$-modules, there is
a long exact sequence
\[
  0 \to T^0(N'') \to T^0(N) \to T^0(N')
\]
\[
  \to T^1(N'') \to T^1(N) \to T^1(N') \to \cdots .
\]
This proves that $T^\bullet$ is the derived functor of $\Hom_B(-, \bar{M})$, and therefore
coincides with $\Ext^\bullet_B(-, \bar{M})$.

(ii) We first prove $\Tor^A_n(M, B) = 0$ for $n > 0$. For $n > 1$ this follows
from $\projdim_A B = 1$. For $n = 1$, consider the long exact sequence
$0 \to \Tor^A_1(M, B) \to M \xrightarrow{x} M \to M \otimes_A B \to 0$ associated with the short
exact sequence $0 \to A \xrightarrow{x} A \to B \to 0$. Since $x$ is $M$-regular we have
$\Tor^A_1(M, B) = 0$.

Now let $L_\bullet \to M \to 0$ be a free resolution of the $A$-module $M$. Then
$L_\bullet \otimes_A B \to M \otimes_A B \to 0$ is exact by what we have just proved, so that
$L_\bullet \otimes_A B$ is a free resolution of the $B$-module $M \otimes_A B = \bar{M}$. Then $\Ext^n_A(M, N)
= H^n(\Hom_A(L_\bullet, N)) = H^n(\Hom_B(L_\bullet \otimes_A B, N)) = \Ext^n_B(\bar{M}, N)$ by Formula 9
of Appendix A.

(iii) Using the same notation as above, we have $\Tor^A_n(M, N) =
H_n(L_\bullet \otimes_A N) = H_n((L_\bullet \otimes_A B) \otimes_B N) = \Tor^B_n(\bar{M}, N)$.
\end{proof}

\begin{lemma}\label{lem:18.3}
Let $(A, m, k)$ be a Noetherian local ring, $M$ a finite $A$-module, and
$P \in \Spec A$; suppose that $\htop (m/P) = 1$. Then
\[
  \Ext^{r+1}_A(k, M) = 0 \Rightarrow \Ext^r_A(A/P, M) = 0.
\]
\end{lemma}

\begin{proof}
Choose $x \in m - P$; then $0 \to A/P \xrightarrow{x} A/P \to A/(P + Ax) \to 0$ is
an exact sequence, and $P + Ax$ is an $m$-primary ideal, so that if we let
$N = A/(P + Ax)$, there exists a chain of submodules of $N$
\[
  N = N_0 \supset N_1 \supset \cdots \supset N_s = 0 \quad \text{such that } N_i/N_{i+1} \cong k.
\]
Hence from $\Ext^{r+1}_A(k, M) = 0$ we get $\Ext^{r+1}_A(A/(P + Ax), M) = 0$, and
\[
  \Ext^r_A(A/P, M) \xrightarrow{x} \Ext^r_A(A/P, M) \to 0
\]
is exact, so that by NAK $\Ext^r_A(A/P, M) = 0$.
\end{proof}

\begin{lemma}\label{lem:18.4}
Let $(A, m, k)$ be a Noetherian local ring, $M$ a finite $A$-module, and
$P \in \Spec A$; suppose that $\htop (m/P) = d$. Then
\[
  \Ext^{r+d}_A(k, M) = 0 \Rightarrow \Ext^r_{A_P}(\kappa(P), M_P) = 0.
\]
\end{lemma}

\begin{proof}
Let $m = P_0 \supset P_1 \supset \cdots \supset P_d = P$, with $P_i \in \Spec A$ and $\htop (P_i/P_{i+1}) = 1$;
then by the previous lemma, from $\Ext^{r+d}_A(k, M) = 0$ we get
\[
  \Ext^{r+d-1}_A(A/P_1, M) = 0,
\]
and localising at $P_1$ we get
\[
  \Ext^{r+d-1}_{A_{P_1}}(\kappa(P_1), M_{P_1}) = 0.
\]
Proceeding in the same way gives the result.
\end{proof}

\begin{theorem}\label{thm:18.1}
Let $(A, m, k)$ be an $n$-dimensional Noetherian local ring.
Then the following conditions are equivalent:
\begin{enumerate}[(1)]
  \item $\injdim A < \infty$;
  \item $\injdim A = n$;
  \item $\Ext^i_A(k, A) = 0$ for $i \neq n$ and $\cong k$ for $i = n$;
  \item $\Ext^i_A(k, A) = 0$ for some $i > n$;
  \item $\Ext^i_A(k, A) = 0$ for $i < n$ and $\cong k$ for $i = n$;
  \item[($4'$)] $A$ is a CM ring and $\Ext^n_A(k, A) \cong k$;
  \item[($5$)] $A$ is a CM ring, and every parameter ideal of $A$ is irreducible;
  \item[($5'$)] $A$ is a CM ring and there exists an irreducible parameter ideal.
\end{enumerate}
Recall that an ideal $I$ is \emph{irreducible} if $I = J \cap J'$ implies either $I = J$ or
$I = J'$ (see \S\ref{sec:6}).
\end{theorem}

\begin{definition}
A Noetherian local ring for which the above equivalent
conditions hold is said to be \emph{Gorenstein}.\index{Gorenstein ring}
\end{definition}

\begin{proof}[Proof of $(1) \Rightarrow (1')$]
Set $\injdim A = r$. If $P$ is a minimal prime ideal of $A$
such that $\htop (m/P) = \dim A = n$ then $PA_P \in \Ass(A_P)$, so that
$\Hom_{A_P}(\kappa(P), A_P) \neq 0$; hence, by Lemma~\hyperref[lem:18.4]{4}, $\Ext^n_A(k, A) \neq 0$, therefore $r \ge n$.
If $r = 0$ this means that $n = 0$, and we are done. If $r > 0$, set $\Ext^r_A(-, A) =
T$; then this is a right-exact contravariant functor, and by Lemma~\hyperref[lem:18.1]{1},
there is a prime ideal $P$ such that $T(A/P) \neq 0$. Now if $P \neq m$ and we take
$x \in m - P$, the exact sequence
\[
  0 \to A/P \xrightarrow{x} A/P
\]
leads to an exact sequence
\[
  T(A/P) \xrightarrow{x} T(A/P) \to 0;
\]
but then by NAK, $T(A/P) = 0$, which is a contradiction. Thus $P = m$, and
so $T(k) \neq 0$. We have $m \notin \Ass(A)$, since otherwise there would exist an
exact sequence $0 \to k \to A$, and hence an exact sequence
\[
  T(A) = \Ext^r_A(A, A) = 0 \to T(k) \to 0,
\]
which is a contradiction. Hence $m$ contains an $A$-regular element $x$. If we set
$B = A/xA$ then by Lemma~\hyperref[lem:18.2]{2}, $\Ext^n_B(N, B) = \Ext^{n+1}_A(N, A)$ for every $B$-
module $N$, so that $\injdim B = r - 1$. By induction on $r$ we have $r - 1 =
\dim B = n - 1$, and hence $r = n$.
\end{proof}

\begin{proof}[Proof of $(1') \Rightarrow (2)$]
When $n = 0$ we have $m \in \Ass(A)$, so there exists an exact
sequence $0 \to k \to A$, and since $\injdim A = 0$,
\[
  A = \Hom(A, A) \to \Hom(k, A) \to 0
\]
is exact. Therefore $\Hom(k, A)$ is generated by one element. But
$\Hom (k, A) \neq 0$, so that we must have $\Hom (k, A) \cong k$. By assumption, $A$ is
an injective module, so that $\Ext^i_A(k, A) = 0$ for $i > 0$; thus we are done in the
case $n = 0$. If $n > 0$ then, as we have seen above, $m$ contains an $A$-regular
element $x$, and if we set $B = A/xA$ then $\dim B = \injdim B = n - 1$, so that by
Lemma~\hyperref[lem:18.2]{2} and induction on $n$ we have
\[
  \Ext^i_A(k, A) = \Ext^{i-1}_B(k, B) =
  \begin{cases}
    0 & \text{if } 0 < i \neq n \\
    k & \text{if } i = n
  \end{cases}
\]
and
\[
  \Hom_A(k, A) = 0. \qquad\blacksquare
\]
\end{proof}

\begin{proof}[Proof of $(3) \Rightarrow (1)$]
We use induction on $n$. Assume that for some $i > n$ we have
$\Ext^i_A(k, A) = 0$. If $n = 0$ then $m$ is the unique prime ideal of $A$, so that by
Lemma~\hyperref[lem:18.1]{1}, $\injdim A \le i < \infty$. If $n > 0$ let $P$ be a prime ideal distinct
from $m$, and set $d = \htop (m/P)$ and $B = A_P$; then by Lemma~\hyperref[lem:18.4]{4} we have
$\Ext^{i-d}_B(\kappa(P), B) = 0$. Moreover, $\dim B \le n - d < i - d$, so that by induction
$\injdim B < \infty$. Thus for any finite $A$-module $M$ we have
\[
  (\Ext^i_A(M, A))_P = \Ext^i_B(M_P, B) = 0
\]
for $i > n \ge \dim B = \injdim B$). Therefore, setting $T(M) = \Ext^i_A(M, A)$,
we get $\Supp(T(M)) = \{m\}$, and since $T(M)$ is a finite $A$-module,
$T(M) \neq 0$ implies that $T(M) \otimes_A k \neq 0$. Using this, we now prove that $T(A/P) = 0$ for every prime
ideal $P$. If $T(A/P) \neq 0$ for some $P$, choose a maximal $P$ with this property. By
assumption $T(k) = 0$, so that $P \neq m$, so that we can take $x \in m - P$ and form
\[
  0 \to A/P \xrightarrow{x} A/P \to A/(P + Ax) \to 0.
\]
We write $A/(P + Ax) = M_0 \supset M_1 \supset \cdots \supset M_s = 0$ with $M_l/M_{l+1} \cong A/P_l$;
then $P_l$ is strictly bigger than $P$, so that $T(A/(P + Ax)) = 0$. Therefore
\[
  T(A/P) \xrightarrow{x} T(A/P) \to 0
\]
so that multiplication by $x$ in $T(A/P)$ is injective; but since
$l(T(A/P)) < \infty$, injective implies surjective. Hence by NAK, $T(A/P) = 0$,
which is a contradiction. Therefore $T(A/P) = 0$ for every $P \in \Spec A$, so that
$\injdim A < \infty$ by Lemma~\hyperref[lem:18.1]{1}.

Thus we have proved that $(1), (1'), (2)$ and $(3)$ are equivalent. Now we
prove the equivalence of $(2), (4), (4'), (5)$ and $(5')$.

$(4) \Rightarrow (2)$ is obvious. $(4) \Leftrightarrow (4')$ comes at once from the fact that $A$ is CM if and
only if $\Ext^i_A(k, A) = 0$ for all $i < n$ (the implication $(2) \Rightarrow (3)$ of
Theorem~\ref{thm:16.6}).

$(4') \Rightarrow (5)$. A system of parameters $x_1, \ldots, x_n$ in a CM ring $A$ is an $A$-
sequence, so that setting $B = A/\sum x_iA$, we have
\[
  \Hom_A(k, B) \cong \Ext^n_A(k, A) \cong k.
\]
Then $B$ is an Artinian ring, and any minimal non-zero ideal of $B$ is
isomorphic to $k$, so that the above formula says that $B$ has just one such
ideal, say $I_0$. If $I_1$ and $I_2$ are any non-zero ideals of $B$ then both of
them contain $I_0$, so that $I_1 \cap I_2 \neq (0)$. Lifting this up to $A$, this means
that $(x_1, \ldots, x_n)$ is an irreducible ideal.

$(5') \Rightarrow (2)$. If $A$ is CM we already have $\Ext^i_A(k, A) = 0$ for $i < n$. If $q$ is
an irreducible parameter ideal and we set $B = A/q$ then, in the same way as
above,
\[
  \Ext^{n+1}_A(k, A) \cong \Ext^1_B(k, B),
\]
so that it is enough to prove that in an Artinian ring $B$, $(0)$ is irreducible
implies that
\[
  \Hom_B(k, B) \cong k \quad \text{and} \quad \Ext^i_B(k, B) = 0 \text{ for } i > 0.
\]
The statement for Hom is easy: first of all, $B$ is Artinian, so that
$\Hom_B(k, B) \neq 0$; for non-zero $f, g \in \Hom_B(k, B)$ we must have $f(k) = g(k)$,
since otherwise $f(k) \cap g(k) = (0)$, which contradicts the irreducibility of $(0)$.
Hence $f(1) = g(\alpha)$ for some $\alpha \in k$, and $f = \alpha g$, so that $\Hom_B(k, B) = k$.

Now consider the $\Ext^i_B(k, B)$. Choose a chain of ideals $(0) = N_0 \subset
N_1 \subset \cdots \subset N_r = B$ such that $N_i/N_{i-1} \cong k$, and consider the exact sequences
\[
  0 \to N_1 \to N_2 \to k \to 0
\]
\[
  0 \to N_2 \to N_3 \to k \to 0
\]
\[
  \vdots
\]
\[
  0 \to N_{r-1} \to B \to k \to 0.
\]
From the long exact sequence
\[
  0 \to \Hom_B(k, B) \to \Hom_B(N_{i+1}, B) \to \Hom_B(N_i, B) \xrightarrow{\delta_i}
  \Ext^1_B(k, B) \to \cdots
\]
and an easy induction (using $N_1 \cong k$ and $\Hom_B(k, B) \cong k$), we get that
$l(\Hom_B(N_i, B)) \le i$, with equality holding if and only if $\delta_1, \ldots, \delta_{i-1}$ are all
zero. However,
\[
  l(\Hom_B(N_r, B)) = l(\Hom_B(B, B)) = l(B) = r,
\]
so that we must have $\delta_1 = \cdots = \delta_{r-1} = 0$. Then from
\[
  0 \to N_{r-1} \to B \to k \to 0
\]
we get the exact sequence
\[
  0 \to \Ext^1_B(k, B) \to \Ext^1_B(B, B) = 0,
\]
and therefore $\Ext^1_B(k, B) = 0$. Now from Lemma~\hyperref[lem:18.1]{1}, $B$ is an injective $B$-
module, so that $\Ext^i_B(k, B) = 0$ for all $i > 0$.
\end{proof}

\begin{lemma}\label{lem:18.5}
Let $A$ be a Noetherian ring, $S \subset A$ a multiplicative set, and $I$ an
injective $A$-module; then $I_S$ is an injective $A_S$-module.
\end{lemma}

\begin{proof}
Every ideal of $A_S$ is the localisation $\mathfrak{a}_S$ of an ideal $\mathfrak{a}$ of $A$. From
$0 \to \mathfrak{a} \to A$ we get the exact sequence $\Hom_A(A, I) \to \Hom_A(\mathfrak{a}, I) \to 0$,
and, since $\mathfrak{a}$ is finitely generated
\[
  \Hom_{A_S}(A_S, I_S) \to \Hom_{A_S}(\mathfrak{a}_S, I_S) \to 0
\]
is exact. This proves that $I_S$ is an injective $A_S$-module.
\end{proof}

\begin{theorem}\label{thm:18.2}
If $A$ is a Gorenstein local ring and $P \in \Spec A$ then $A_P$ is also
Gorenstein.
\end{theorem}

\begin{proof}
If $0 \to A \to I^0 \to I^1 \to \cdots \to I^n \to 0$ is an injective reso-
lution of $A$, then localising at $P$ and applying the previous lemma gives an
injective resolution of $A_P$, so that $\injdim A_P \le n$.
\end{proof}

\begin{definition}
A Noetherian ring $A$ is \emph{Gorenstein} if its localisation at every
maximal ideal is a Gorenstein local ring. (By the previous theorem, it then
follows that $A_P$ is Gorenstein for every $P \in \Spec A$.)
\end{definition}

\begin{theorem}\label{thm:18.3}
Let $A$ be a Noetherian local ring and $\hat{A}$ its completion.
Then $A$ is Gorenstein $\iff \hat{A}$ is Gorenstein.
\end{theorem}

\begin{proof}
We have $\dim A = \dim \hat{A}$, and since $\hat{A}$ is faithfully flat over $A$,
$\Ext^i_A(k, A) \otimes_A \hat{A} = \Ext^i_{\hat{A}}(k, \hat{A})$, so that we only need to use condition (3)
of Theorem 1.
\end{proof}

Closely related to the theory of Gorenstein rings is Matlis' theory
of injective modules over Noetherian rings. We now discuss the main
results of Matlis~\cite{Matl1}.

Let $A$ be a Noetherian ring, and $E$ an injective $A$-module. If $E$ is a
submodule of an $A$-module $M$ then since we can extend the identity map
$1_E$ to a linear map $f\colon M \to E$, we have $M = E \oplus F$ (with $F = \Ker f$).
An $A$-module $N$ is \emph{indecomposable}\index{indecomposable module} if $N$ cannot be written as a direct
sum of two submodules. We write $E_A(N)$ or $E(N)$ for the \emph{injective hull}\index{injective hull}
of an $A$-module $E$ (see Appendix B).

\begin{theorem}\label{thm:18.4}
Let $A$ be a Noetherian ring and $P, Q \in \Spec A$.
\begin{enumerate}[(i)]
  \item $E(A/P)$ is indecomposable.
  \item Any indecomposable injective $A$-module is of the form $E(A/P)$ for
        some $P \in \Spec A$.
  \item If $x \in A - P$ then multiplication by $x$ induces an automorphism of $E(A/P)$.
  \item $P \nsubseteq Q \iff E(A/P) \to E(A/Q)$ is injective.
  \item For any $\xi \in E(A/P)$ there exists a positive integer $\nu$ (depending on $\xi$)
        such that $P^\nu\xi = 0$.
  \item If $Q \subset P$ then $E(A/Q)$ is an $A_P$-module, and is an injective hull of
        $(A/Q)_P = A_P/QA_P$, that is
        \[
          E_A(A/Q) = E_{A_P}(A_P/QA_P).
        \]
\end{enumerate}
\end{theorem}

\begin{proof}
(i) If $I_1$ and $I_2$ are non-zero ideals of $A/P$ then $0 \neq I_1I_2 \subset I_1 \cap I_2$.
Since $E(A/P)$ is an essential extension of $A/P$ (see Appendix B), so that for
any two non-zero submodules $N_1, N_2$ of $E(A/P)$ we have $N_i \cap (A/P) \neq 0$, so
that
\[
  N_1 \cap N_2 \supset (N_1 \cap A/P) \cap (N_2 \cap A/P) \neq 0.
\]
(ii) Let $N \neq 0$ be an indecomposable injective $A$-module and choose
$P \in \Ass(N)$. Then $A/P$ can be embedded into $N$, and so $E(A/P)$ can also; but
an injective submodule is always a direct summand, and since $N$ is
indecomposable, $N = E(A/P)$.

(iii) Write $\varphi$ for multiplication by $x$ in $E(A/P)$; then $\Ker(\varphi) \cap (A/P) = 0$, so
that $\Ker(\varphi) = 0$, and $\Im \varphi$ is isomorphic to $E(A/P)$. Hence $\Im \varphi$ is injective,
and is therefore a direct summand of $E(A/P)$, so that by (i), $\Im \varphi = E(A/P)$.

(iv) If $P \nsubseteq Q$ and $x \in P - Q$ then multiplication by $x$ is injective in
$E(A/Q)$ but not in $E(A/P)$.

(v) By the proof of (ii) together with (iv), $\Ass(E(A/P)) = \{P\}$, so that the
submodule $A\xi \cong A/\ann(\xi)$ also has $\Ass(A\xi) = \{P\}$. Hence $\ann(\xi)$ is a
$P$-primary ideal.

(vi) By (iii), we can view $E(A/Q)$ as an $A_P$-module; hence it contains
$(A/Q)_P$. Since $E(A/Q)$ is an essential extension of $A/Q$ and $A/Q \subset (A/Q)_P \subset
E(A/Q)$, it is also an essential extension of $(A/Q)_P$. For $A_P$-modules $M$ and $N$,
any $A$-linear map $M \to N$ is also $A_P$-linear, and of course conversely,
so that for an $A_P$-module, being injective as an $A_P$-module is the same as
being injective as an $A$-module. Thus $E(A/Q)$ is an injective hull of the
$A_P$-module $(A/Q)_P$.
\end{proof}

\begin{example}\label{ex:18.1}
If $A$ is an integral domain and $K$ its field of fractions, $K = E(A)$
(prove this!).
\end{example}

\begin{example}\label{ex:18.2}
If $A$ is a DVR with uniformising element $x$ and field of fractions
$K$, and $k = A/xA$, then $E(k) = K/A$. Indeed, if $I$ is a non-zero ideal of $A$
we can write $I = x^rA$, and if $f\colon I \to K/A$ is a given map, let $f(x^r) = \alpha
\bmod A$ for some $\alpha \in K$; then $f$ can be extended to a map $\bar{f}\colon A \to K/A$
by setting $\bar{f}(1) = (\alpha/x^r) \bmod A$. Therefore $K/A$ is injective. We have
$(x^{-1}A)/A \cong A/xA = k$, and it is easy to see that $K/A$ is an essential extension
of $x^{-1}A/A$. Thus $K/A$ can be thought of as $E(k)$.
\end{example}

\begin{theorem}\label{thm:18.5}
We consider modules over a Noetherian ring $A$.
\begin{enumerate}[(i)]
  \item A direct sum of any number of injective modules is injective.
  \item Every injective module is a direct sum of indecomposable injective
        modules.
  \item The direct sum decomposition in (ii) is unique, in the sense that if
        \[
          M = \bigoplus M_i \quad \text{(with indecomposable } M_i\text{)}
        \]
        then for any $P \in \Spec A$, the sum $M(P)$ of all the $M_i$ isomorphic to $E(A/P)$
        depends only on $M$ and $P$, and not on the decomposition $M = \bigoplus M_i$.
        Moreover, the number of $M_i$ isomorphic to $E(A/P)$ is equal to
        \[
          \dim_{\kappa(P)} \Hom_{A_P}(\kappa(P), M_P), \quad \text{(where } \kappa(P) = A_P/PA_P),
        \]
        so that this also is independent of the decomposition.
\end{enumerate}
\end{theorem}

\begin{proof}
(i) Let $M_\lambda$ for $\lambda \in \Lambda$ be injective modules. It is enough to prove that
for an ideal $I$ of $A$, any linear map $\varphi\colon I \to \bigoplus M_\lambda$ can be extended to a linear
map $\bar{\varphi}\colon A \to \bigoplus M_\lambda$. Since $I$ is finitely generated, $\varphi(I)$ is contained
in a finite direct sum $M_{\lambda_1} \oplus \cdots \oplus M_{\lambda_n}$ of the $M_\lambda$. If $\varphi_\mu$ is the $\mu$th component
of $\varphi(\mathfrak{a})$, then $\varphi_\mu\colon I \to M_\mu$ extends
to $\bar{\varphi}_\mu\colon A \to M_\mu$, and we just set $\bar{\varphi}(1) = \bar{\varphi}_1(1) + \cdots + \bar{\varphi}_n(1)$.

(ii) Say that a family $\mathcal{F} = \{E_\lambda\}$ of indecomposable injective submodules
$E_\lambda$ of $M$ is \emph{free} if the sum in $M$ of the $E_\lambda$ is direct, that is if, for any finite
number $E_{\lambda_1}, \ldots, E_{\lambda_n}$ of them,
\[
  E_{\lambda_n} \cap (E_{\lambda_1} + \cdots + E_{\lambda_{n-1}}) = 0.
\]
Let $\mathfrak{M}$ be the set of all free families $\mathcal{F}$, ordered by inclusion. Then by
Zorn's lemma $\mathfrak{M}$ has a maximal element, say $\mathcal{F}_0$. Write $N = \sum_{E \in \mathcal{F}_0} E$; then
by (i), $N$ is injective, hence a direct summand of $M$, and $M = N \oplus N'$. If
$N' \neq 0$ then since it is a direct summand of $M$ it must be injective, and
for $P \in \Ass(N')$, the proof of Theorem 4, (ii), shows that $N'$ contains a
direct summand $E'$ isomorphic to $E(A/P)$. Thus $\mathcal{F}_0 \cup \{E'\}$ is a free family,
contradicting the maximality of $\mathcal{F}_0$. Hence $N' = 0$ and $M = N$.

(iii) If we can show that $M(P)$ has the property that every submodule
$E$ of $M$ isomorphic to $E(A/P)$ is contained in $M(P)$, then $M(P)$ is the
submodule of $M$ generated by all such $E$, and therefore is determined by
$M$ and $P$ only. To prove this, take any $\xi \in E$; we can write $\xi = \xi_1 + \cdots +
\xi_r$ with $\xi_i \in M(P_i)$, where $P_1, \ldots, P_r$ are distinct prime ideals and $P = P_1$.
Setting $\eta_1 = \xi - \xi_1$, and $\eta_i = \xi_i$ for $2 \le i \le r$ we have $\eta_1 + \cdots + \eta_r = 0$, with
$\eta_1 \in M(P_1) + E$ and $\eta_i \in M(P_i)$ for $i \ge 2$. We need only prove that in this case
each $\eta_i = 0$. Suppose that $P_1$ is minimal among $P_1, \ldots, P_r$; then for any $m$
we have $(P_1 \cdots P_r)^m \nsubseteq P_1$, so that taking $a \in (P_1 \cdots P_{r-1})^m - P_1$ and $m$
large enough, we get $a\eta_1 = \cdots = a\eta_{r-1} = 0$. Then also $\eta_r = 0$, but
multiplication by $a$ is an automorphism of $M(P_r)$, so that $\eta_r = 0$. By
induction on $r$ we get $\eta_i = 0$ for all $i$.

We now prove that if $M(P) = M_1 \oplus \cdots \oplus M_s$ with $M_i \cong E(A/P)$ then
\[
  s = \dim_{\kappa(P)} \Hom_{A_P}(\kappa(P), M_P).
\]
(We are writing this as if $s$ were finite, but, as one can see from the proof
below, the same works for any cardinal number.) By Theorem 4, (vi), both
sides of $M(P) = M_1 \oplus \cdots \oplus M_s$ are $A_P$-modules, and $M_i = E(\kappa(P))$.
Moreover, by Theorem 4, (v), $E(A/Q)_P = 0$ if $Q \nsubseteq P$, so that
\[
  M_P = \bigoplus_{Q \subset P} M(Q)_P = \bigoplus_{Q \subset P} M(Q).
\]
Hence we can replace $A$ by $A_P$, and assume that $A$ is a local ring with $P$
its maximal ideal; set $k = \kappa(P)$. If $Q \neq P$ then any $x \in P - Q$ gives an
automorphism of $M(Q)$, but $x \cdot k = 0$, so that $\Hom_A(k, M(Q)) = 0$. Hence
$\Hom_A(k, M) = \Hom_A(k, M(P))$, so that there is no loss of generality in
assuming that $M = M(P)$. For any $A$-module $N$, we can identify
$\Hom_A(k, N)$ with the submodule $\{\xi \in N \mid P\xi = 0\}$, but since $E(k)$ is an
essential extension of $k$ we must have $\dim_k \Hom_A(k, E(k)) = 1$, so that if
$M = M_1 \oplus \cdots \oplus M_s$ with $M_i \cong E(k)$ then $s = \dim_k \Hom_A(k, M)$.
\end{proof}

\begin{theorem}\label{thm:18.6}
Let $(A, m, k)$ be a Noetherian local ring, and $E = E_A(k)$ the
injective hull of $k$. For each $A$-module $M$ set $M' = \Hom_A(M, E)$.
\begin{enumerate}[(i)]
  \item If $M$ is an $A$-module and $0 \neq x \in M$, then there exists $\varphi \in M'$ such that
        $\varphi(x) \neq 0$. In other words the canonical map $\theta\colon M \to M''$ defined by $\theta(x)(\varphi)
        = \varphi(x)$ for $x \in M$ and $\varphi \in M'$ is injective.
  \item If $M$ is an $A$-module of finite length, then $l(M) = l(M')$ and the
        canonical map $M \to M''$ is an isomorphism.
  \item Let $\hat{A}$ be the completion of $A$; then $E$ is also an $\hat{A}$-module, and
        is an injective hull of $k$ as $\hat{A}$-module.
  \item $\Hom_A(E, E) = \Hom_{\hat{A}}(E, E) = \hat{A}$. In other words, each endomor-
        phism of the $A$-module $E$ is multiplication by a unique element of $\hat{A}$.
  \item $E$ is Artinian as an $A$-module and also as an $\hat{A}$-module. Assume
        now that $A$ is complete, and write $\mathcal{N}$ (resp.\ $\mathcal{A}$) for the category of
        Noetherian (respectively Artinian) $A$-modules. Then if $M \in \mathcal{A}$ we have
        $M' \in \mathcal{N}$ and $M \cong M''$; if $M \in \mathcal{N}$ we have $M' \in \mathcal{A}$ and $M \cong M''$.
\end{enumerate}
\end{theorem}

\begin{proof}
(i) Let $f\colon Ax \to E$ be the composite of the canonical maps
$Ax \cong A/\ann(x)$, $A/\ann(x) \to A/m = k$ and $k \to E$. Then $f(x) \neq 0$. Since $E$ is
injective we can extend $f$ to $\varphi\colon M \to E$.

(ii) If $l(M) = n < \infty$ then $M$ has a submodule $M_1$ of length $n - 1$, and
$0 \to M_1 \to M \to k \to 0$ is exact, so that $0 \to k' \to M' \to M_1' \to 0$ is exact.
However
\[
  k' = \Hom_A(k, E) = \Hom_A(k, k) \cong k,
\]
so that by induction on $n$ we get $l(M) = n = l(M')$. The canonical map
$M \to M''$ is injective by (i), and $l(M) = l(M') = l(M'')$, hence it must be an
isomorphism.

(iii) Each element of $E$ is annihilated by some power of $m$, so that the
canonical map $E \to E \otimes_A \hat{A}$ is surjective. However, since $\hat{A}$ is faithfully
flat over $A$ it is also injective, so that $E \cong E \otimes_A \hat{A}$, and we can view $E$ as
an $\hat{A}$-module. Let $F$ be the injective hull of $E$ as an $\hat{A}$-module. Then $F$
is also the $A$-injective hull of $k$, so that every element of $F$ is annihilated
by some power of $m\hat{A}$. As an $A$-module $F$ splits into a direct sum of $E$
and an $\hat{A}$-module $C$. If $x \in C$, and if $m^\nu Ax = 0$, then for each $\hat{a} \in \hat{A}$ we
can find $a \in A$ such that $\hat{a} \equiv a \bmod m^\nu\hat{A}$ and hence $\hat{a}x = ax \in C$.
Therefore $C$ is an $A$-module. But $F$ is indecomposable as an $\hat{A}$-module.
Hence $C = 0$ and $E = F$.

(iv) For $\nu > 0$ set $E_\nu = \{x \in E \mid m^\nu x = 0\}$. Then we have $(A/m^\nu)' \cong
\Hom_A(A/m^\nu, E) \cong E_\nu$, and $\Hom_A(E_\nu, E_\nu) = \Hom_A(E_\nu, E) = E'_\nu = (A/m^\nu)'' =
A/m^\nu$. Now $E_1 \subset E_2 \subset \cdots$ and $E = \bigcup_\nu E_\nu$ by Theorem 4, (v), hence $E =
\varinjlim E_\nu$. Therefore $\Hom_A(E, E) = \Hom_A(\varinjlim E_\nu, E) = \varprojlim \Hom_A(E_\nu, E) =
\varprojlim A/m^\nu = \hat{A}$.

(v) If an $A$-module $M$ is Artinian and $x \in M$, then $Ax \cong A/\ann(x)$ is also
Artinian and consequently $m^\nu \subset \ann(x)$ for some $\nu$. Therefore $M$ can be
viewed as an $\hat{A}$-module, and its $\hat{A}$-submodules are precisely its $A$-
submodules. It is also clear that if an $\hat{A}$-module $M$ is Artinian then we have
the same conclusion. Therefore to prove (v) we may assume that $A$ is
complete.

If $M$ is a submodule of $E$ set $M^\perp = \{a \in A \mid aM = 0\}$. If $I$ is an ideal of $A$ set
$I^\perp = \{x \in E \mid Ix = 0\}$. Then clearly $M^{\perp\perp} \supset M$. If $x \in E - M$ there exist
$\varphi \in (E/M)'$ satisfying $\varphi(x \bmod M) \neq 0$ by (i), and if we identify $E' =
\Hom_A(E, E)$ with $A$ then $(E/M)'$ is identified with $M^\perp$. Thus $\varphi(x \bmod M) =
ax$ for some $a \in M^\perp$, and $x \notin M^{\perp\perp}$. Therefore $M^{\perp\perp} = M$. Similarly, if
$a \in \hat{A} - I$ then there exists $\varphi \in (A/I)'$ such that $\varphi(a \bmod I) \neq 0$, and
$(A/I)'$ is identified with the submodule $I^\perp$ of $E = A'$. Thus, setting $x =
\varphi(1 \bmod I)$ we have $x \in I^\perp$ and $ax = \varphi(a \bmod I) \neq 0$. This proves $a \notin I^{\perp\perp}$, so
that $I = I^{\perp\perp}$. Thus $M \mapsto M^\perp$ is an order-reversing bijection from the set of
submodules of $E$ onto the set of ideals of $A$. Since $A$ is Noetherian, it follows
that $E$ is Artinian. By Theorem~\ref{thm:3.1} finite direct sums $E^n$ of $E$ are also
Artinian for all $n > 0$.

If $M \in \mathcal{N}$ then there is a surjection $A^n \to M$ for some $n$, and so
there is an injection $M' \hookrightarrow (A^n)' = E^n$. Hence $M'$ is Artinian. On the
other hand, if $M \in \mathcal{A}$ there is an injection $M \to E^n$ for some $n$. This can be
seen as follows: consider all linear maps $M \to E^n$, where $n$ is not fixed,
and take one $\varphi\colon M \to E^n$ whose kernel is minimal among the kernels of
those maps. Then, using (i) we can easily see that $\Ker(\varphi) = 0$. Now, from
$0 \to M \to E^n$ we have $(E^n)' = A^n \to M' \to 0$ exact, hence $M' \in \mathcal{N}$. Now the
assertion $M \cong M''$ for $M \in \mathcal{N}$ or $\mathcal{A}$ can easily be checked using (iv) if $M = A$
or $E$, and the general case follows from this and from (i).
\end{proof}

\begin{lemma}\label{lem:18.6}
Let $A$ be a Noetherian ring, $S \subset A$ a multiplicative set, $M$ an $A$-
module and $N \subset M$ a submodule. Assume that $M$ is an essential extension
of $N$; then $M_S$ is an essential extension of $N_S$.
\end{lemma}

\begin{proof}
For $\xi \in M$ we write $\xi_s = \xi/1 \in M_S$; then any element of $M_S$ can be written
$u^{-1}\xi_s$ (with $u$ a unit of $A$, and $\xi \in M$), so that it is enough to show that
for any non-zero $\xi_s$ we have $N_S \cap A_S\xi_s \neq 0$. Suppose that $\ann(\xi_s)$ is a
maximal element of the set of ideals $\{\ann(t\xi_s) \mid t \in S\}$; then if we set $\eta = \xi_s$,
we have $\eta_s = \eta$, and hence $\eta \neq 0$. Now let $\mathfrak{b} = \{b \in A \mid b\eta \in N\}$; by
assumption,
\[
  \mathfrak{b}\eta = A\eta \cap N \neq 0.
\]
Suppose that $\mathfrak{b} = (b_1, \ldots, b_r)$; if $b_1\eta_s = \cdots = b_r\eta_s = 0$ then there is a $t \in S$
such that $tb_i\eta = 0$ for all $i$. Then $t\mathfrak{b}\eta = 0$, but by choice of $\eta$ we have
$\ann(\eta) = \ann(t\eta)$, so that $\mathfrak{b}\eta = 0$, which is a contradiction. Thus $b_i\eta_s \neq 0$
for some $i$, and
\[
  b_i\eta_s \in A_S\eta_s \cap N_S = A_S\xi_s \cap N_S,
\]
as required.
\end{proof}

By Lemmas 5 and 6, if $M$ is an injective hull of $N$ then the $A_S$-module
$M_S$ is an injective hull of $N_S$. Hence if $0 \to M \to I^0 \to I^1 \to \cdots$ is
a minimal injective resolution of an $A$-module $M$, then $0 \to M_S \to I^0_S \to
I^1_S \to \cdots$ is a minimal injective resolution of the $A_S$-module $M_S$. The
$I^i$ are determined uniquely up to isomorphism by $M$. We can therefore
define $\mu_i(P, M)$ to be the number of summands isomorphic to $E(A/P)$
appearing in a decomposition of $I^i$ as a direct sum of indecomposable
modules. We can write symbolically
\[
  I^i = \bigoplus_{P \in \Spec A} \mu_i(P, M)E(A/P).
\]
From what we have just proved, for a multiplicative set $S \subset A$,
\[
  \mu_i(P, M) = \mu_i(PA_S, M_S) \quad \text{if } P \cap S = \emptyset.
\]

\begin{theorem}\label{thm:18.7}
Let $A$ be a Noetherian ring, $M$ an $A$-module, and
$P \in \Spec A$. Then
\[
  \mu_i(P, M) = \dim_{\kappa(P)} \Ext^i_{A_P}(\kappa(P), M_P) = \dim_{\kappa(P)}(\Ext^i_A(A/P, M))_P.
\]
In particular, if $M$ is a finite $A$-module then $\mu_i(P, M) < \infty$.
\end{theorem}

\begin{proof}
Replacing $A$ and $M$ by $A_P$ and $M_P$ we can assume that $(A, P, k)$ is
a local ring. Let $0 \to M \to I^0 \to I^1 \to \cdots$ be a minimal injective
resolution of $M$, so that $\Ext^i_A(k, M)$ is obtained as the homology of the
complex
\[
  \cdots \to \Hom_A(k, I^{i-1}) \to \Hom_A(k, I^i) \to \Hom_A(k, I^{i+1})
  \to \cdots
\]
We can identify $\Hom_A(k, I^i)$ with the submodule $T^i = \{x \in I^i \mid Px = 0\} \subset
I^i$. By construction of the minimal injective resolution, $I^i$ is an essential
extension of $d(I^{i-1})$, so that for $x \in T^i$ the submodule $Ax \cong k$ intersects
$d(I^{i-1})$, and $x \in d(I^{i-1})$. Therefore, $T^i \subset d(I^{i-1})$, and $dT^{i-1} = dT^i = 0$, so
that $\Ext^i_A(k, M) = T^i$. Also,
\[
  \dim_k T^i = \dim_k \Hom_A(k, I^i),
\]
and by Theorem 4, (iii), this is equal to $\mu_i(P, M)$.
\end{proof}

\begin{theorem}\label{thm:18.8}
A necessary and sufficient condition for a ring $A$ to be
Gorenstein is that a minimal injective resolution $0 \to A \to I^0 \to
I^1 \to \cdots$ of $A$ satisfies
\[
  I^i = \bigoplus_{\htop P = i} E(A/P),
\]
or, in other words, $\mu_i(P, A) = \delta_{i,\htop P}$ (the Kronecker $\delta$) for every $P \in \Spec A$.
\end{theorem}

\begin{proof}
By Theorem 7 and condition (2) of Theorem 1 we have
\[
  A_P \text{ is Gorenstein} \iff \mu_i(P, A) = \delta_{i,\htop P}. \qquad\blacksquare
\]
\end{proof}

\begin{theorem}\label{thm:18.9}
Let $(A, m)$ be a Noetherian local ring, and $M$ a finite
$A$-module. Then
\[
  \injdim M < \infty \iff \injdim M = \depth A.
\]
\end{theorem}

\begin{proof}
Suppose that $\injdim M = r < \infty$. If $P$ is a prime ideal distinct from
$m$, choose $x \in m - P$. Then
\[
  0 \to A/P \xrightarrow{x} A/P,
\]
together with the right-exactness of $\Ext^r_A(-, M)$ gives an exact sequence
\[
  \Ext^r_A(A/P, M) \xrightarrow{x} \Ext^r_A(A/P, M) \to 0,
\]
so that by NAK $\Ext^r_A(A/P, M) = 0$. Putting this together with Lemma~\hyperref[lem:18.1]{1},
we get $\Ext^r_A(k, M) \neq 0$. Set $t = \depth A$, and let $x_1, \ldots, x_t \in m$ be a
maximal $A$-sequence; then setting $A/(x_1, \ldots, x_t) = N$ we have $m \in \Ass(N)$.
Hence there exists an exact sequence $0 \to k \to N$, and we must have
$\Ext^t_A(N, M) \neq 0$. The Koszul complex $K(x_1, \ldots, x_t)$ is a projective resolution
of $N = A/(x_1, \ldots, x_t)$, so computing Ext by means of it we see that
\[
  \Ext^t_A(N, M) \cong M/(x_1, \ldots, x_t)M,
\]
and by NAK this is non-zero. Thus $\projdim N = t$, and from $\Ext^r_A(N, M) \neq
0$ we get $t \le r$, whereas from $\Ext^t_A(N, M) \neq 0$ we get $t \ge r$. Hence $t = r$.
\end{proof}

\begin{remark}[the Bass conjecture and the intersection theorem]
Let $(A, m, k)$ be
a Noetherian local ring of dimension $d$. H. Bass~\cite{Ba1} conjectured the
following:
\[
  (B) \quad \text{if there exists a finite } A\text{-module } M (\neq 0) \text{ of finite injective}
\]
\[
  \text{dimension, then } A \text{ is a CM ring.}
\]
According to Theorem 9 this is equivalent to asking that $\injdim M = d$.
The converse of the Bass conjecture is true. Indeed, if $A$ is CM, taking a
maximal $A$-sequence $x_1, \ldots, x_d$, and setting $B = A/(x_1, \ldots, x_d)$ and $E = E(k)$
we have $l_A(B) < \infty$. By Theorem 6, $M = \Hom_A(B, E)$ is also of finite length,
hence is finitely generated. We prove $\injdim_A M \le d$; the Koszul complex
$0 \to A \to A^d \to \cdots \to A^d \to A \to B \to 0$ with respect to $x_1, \ldots, x_d$ provides an $A$-
free resolution of $B$. Now applying the exact functor $\Hom_A(-, E)$ to this
gives the exact sequence $0 \to M \to E \to E^d \to \cdots \to E^d \to E \to 0$. This proves
$\injdim_A M \le d$.

$(B)$ is a special case of the following theorem.
\[
  (C) \quad \text{If } I^\bullet\colon 0 \to I^0 \to \cdots \to I^d \to 0 \text{ is a complex of injective modules such}
\]
\[
  \text{that } H^i(I^\bullet) \text{ is finitely generated for all } i \text{ and } I^\bullet \text{ is not exact, then }
\]
\[
  I^d \neq 0.
\]
Using the theory of dualizing complexes (see~\cite{Rob}) one can prove that $(C)$
is equivalent to the following

\emph{Intersection Theorem.} If $F_\bullet\colon 0 \to F_r \to \cdots \to F_0 \to 0$ is a complex of finitely
generated free modules such that $H_i(F_\bullet)$ has finite length for all $i$ and $F_\bullet$ is
not exact, then $F_r \neq 0$.

$(B)$ was proved by Peskine and Szpiro~\cite{PS1} in some important cases, and
by Hochster~\cite{H} in the equal characteristic case (i.e. when $A$ contains a
field) as a corollary of his existence theorem for the `big CM module', i.e. a
(not necessarily finite) $A$-module with $\depth = \dim A$, see~\cite{H} p. 10 and
p. 70. The intersection theorem was conjectured by Peskine--Szpiro~\cite{PS3} and
by P. Roberts independently. They pointed out that it was also a
consequence of Hochster's theorem. Finally, P. Roberts~\cite{Rob3} settled the
remaining unequal characteristic case of the intersection theorem by using
the advanced technique of algebraic geometry developed by W. Fulton
(\cite{Ful}). Therefore $(B)$, which was known as Bass's conjecture for 24 years,
is now a theorem. Some other conjectures listed in~\cite{H} are still open.
\end{remark}

% ============ 习题 ============
\section*{Exercises to \S\ref{sec:18}}
\emph{Prove the following propositions.}

\begin{exercise}\label{ex:18.1-ex}
Let $(A, m)$ be a Noetherian local ring, $x_1, \ldots, x_n$ an $A$-sequence, and set
$B = A/(x_1, \ldots, x_n)$; then $A$ is Gorenstein $\iff B$ is Gorenstein.
\end{exercise}

\begin{exercise}\label{ex:18.2-ex}
Use the result of Ex.~\ref{ex:18.1-ex} to give another proof of Theorem~\hyperref[thm:18.3]{3}.
\end{exercise}

\begin{exercise}\label{ex:18.3-ex}
If $A$ is Gorenstein then so is the polynomial ring $A[X]$.
\end{exercise}

\begin{exercise}\label{ex:18.4-ex}
Is the ring $R$ of Ex.~\ref{ex:17.2-ex} Gorenstein?
\end{exercise}

\begin{exercise}\label{ex:18.5-ex}
Let $(A, m, k)$ be a local ring; then $E = E_A(k)$ is a faithful $A$-module (that is
$0 \neq a \in A \Rightarrow aE \neq 0$).
\end{exercise}

\begin{exercise}\label{ex:18.6-ex}
Let $(A, m, k)$ be a complete Noetherian local ring and $M$ an $A$-module. If $M$
is a faithful $A$-module and is an essential extension of $k$ then $M \cong E_A(k)$.
\end{exercise}

\begin{exercise}\label{ex:18.7-ex}
Let $k$ be a field, $S = k[X_1, \ldots, X_n]$ and $P = (X_1, \ldots, X_n)$; set $A = S_P$,
$\hat{A} = k[[X_1, \ldots, X_n]]$ and $E = k[X_1^{-1}, \ldots, X_n^{-1}]$. We make $E$ into an
$A$-module by the following multiplication: if $X^\alpha = X_1^{\alpha_1} \cdots X_n^{\alpha_n}$ and $X^{-\beta}
= X_1^{-\beta_1} \cdots X_n^{-\beta_n}$, the product $X^\alpha X^{-\beta}$ is defined to be $X^{\alpha-\beta}$ if $\alpha_i \le \beta_i$ for all $i$
and $0$ otherwise. Then $E = E_A(S/P) = E_{\hat{A}}(k)$. (Use the preceding question;
see also Northcott~\cite{Nort8} for further results. The elements of this $A$-module
$E$ are called \emph{inverse polynomials}; they were defined and used by Macaulay
~\cite{Mac} as early as 1916.)
\end{exercise}

\begin{exercise}\label{ex:18.8-ex}
Let $k$ be a field and $t$ an indeterminate. Consider the subring $A =
k[t^3, t^5, t^7]$ of $k[t]$ and show that $A$ is a one-dimensional CM ring which
is not Gorenstein. How about $k[t^3, t^4, t^5]$ and $k[t^4, t^5, t^6]$?
\end{exercise}

